    \documentclass[lettersize,journal]{IEEEtran}
    
    % ---- Unicode sanitization (IMPORTANT) ----
    \DeclareUnicodeCharacter{202A}{}
    \DeclareUnicodeCharacter{202C}{}
    \DeclareUnicodeCharacter{200E}{}
    \DeclareUnicodeCharacter{200F}{}
    
    % ---- Math & theorem support ----
    \usepackage{amsmath,amssymb,amsfonts}
    \usepackage{amsthm}
    \usepackage{bm}
    \usepackage{algorithm}
    \usepackage{algpseudocode}
    
    \newtheorem{theorem}{Theorem}
    \newtheorem{lemma}{Lemma}
\newtheorem{corollary}{Corollary}
    \newtheorem{definition}{Definition}
    \newtheorem{game}{Game}
    
    % ---- Algorithms ----
    
    % ---- Figures & tables ----
    \usepackage{graphicx}
    \graphicspath{{./}{Experiment/}}
    \usepackage{subcaption}
    \usepackage{booktabs}
    \usepackage{tabularx}
    \usepackage{multirow}
    \usepackage{array}
    \usepackage{adjustbox}
    
    \usepackage{pifont}
    \newcommand{\cmark}{\ding{51}} % checkmark
    \newcommand{\xmark}{\ding{55}} % cross mark
    
    % ---- Text & formatting ----
    \usepackage{cite}
    \usepackage{textcomp}
    \usepackage{url}
    \usepackage{verbatim}
    \usepackage{stfloats}
    \usepackage{enumitem}
    \usepackage{setspace}
    \usepackage{xcolor}
    \usepackage{listings}
    \usepackage{makecell}

    \usepackage{placeins}
    
    % ---- References ----
    \usepackage[capitalize]{cleveref}
    
    \hyphenation{op-tical net-works semi-conduc-tor IEEE-Xplore}
    
    % updated with editorial comments 8/9/2021
    
    \begin{document}
    
    \title{VeRedact-PQ: Scalable Post-Quantum Redaction with Multi-Party Authorization and Verifiable Auditing for Permissioned Blockchains}
    
    %}
    %
    %\thanks{
    %Corresponding author: 
    %Somchart Fugkeaw
    %(e-mail: [somchart@siit.tu.ac.th]
    
    
    %\IEEEpubid{0000--0000/00\$00.00~\copyright~2021 IEEE}
    % Remember, if you use this you must call \IEEEpubidadjcol in the second
    % column for its text to clear the IEEEpubid mark.
    
    \author{
    \IEEEauthorblockN{Rangsimann Sattayarom, Ratchanon Wongwitutai, Natthawat Tungsriworakan, \\
    Tagrid Chongkolrattanapond, Somchart Fugkeaw\\}
    \IEEEauthorblockA{Sirindhorn International Institute of Technology, Thammasat University, Thailand\\
    6622771747@g.siit.tu.ac.th, 
    6622780268@g.siit.tu.ac.th,
    6622772364@g.siit.tu.ac.th,
    6622781175@g.siit.tu.ac.th,
    somchart@siit.tu.ac.th}
    }
    
    \maketitle

\begin{abstract}

Redactable permissioned blockchains enable controlled modification of
committed transactions, but existing approaches face challenges in
distributed authorization, post-quantum security, verifiable
accountability, and scalability under high redaction workloads.
Processing each redaction independently can repeatedly incur
policy-validation, committee, authenticated-state update, and audit
costs. This paper proposes \emph{VeRedact-PQ}, a scalable post-quantum
framework for policy-bound redaction and verifiable auditing in
permissioned blockchains. VeRedact-PQ combines state-bound
post-quantum zero-knowledge validation with threshold committee
authorization and distributed post-quantum chameleon-hash adaptation.
Its Cost- and Overlap-Aware Redaction Aggregation (CORA) identifies
validated requests that can share execution work, while Shared-Frontier
Merkle Redaction (SFMR) exploits this locality through shared Merkle
processing and batch-level chameleon-hash adaptation. Policy-Bound
Batch Redaction Proofs (PBRP) bind request validation and committee
authorization to the finalized state transition, preserving
request-level accountability under shared execution. For auditing, a
sharded authenticated redaction index and State-Aware Verification
Evidence Reuse (SAVER) organize evidence as state-bound dependencies,
allowing common authorization, transition, and checkpoint evidence to
be verified once under matching contexts while selectively escalating
challenged requests to independent post-quantum zero-knowledge
verification. On a 7-validator Hyperledger Besu testbed, VeRedact-PQ
sustains more valid redactions than every baseline from 250 requests/s
while accepting no invalid request, keeps its per-request authorization
cost nearly constant as the committee grows to 64 members, and verifies
$10^4$ audit records about six times faster than a decentralized
chameleon-hash baseline (5.3 vs.\ 33.6~s). A lattice-based batch
authentication scheme remains cheaper for authorization, and an
aggregate audit proof is faster when no record is invalid, but neither
establishes who authorized a redaction or which record is invalid.

\end{abstract}

\begin{IEEEkeywords}
Permissioned Blockchain, Redactable Blockchain, Post-Quantum Security,
Post-Quantum Chameleon Hash, Verifiable Redaction, Privacy-Preserving
Auditing, Scalable Redaction.
\end{IEEEkeywords}

\section{Introduction}

Permissioned blockchains provide a shared and auditable ledger for
multi-organization environments without requiring complete mutual
trust~\cite{ref7}. Their integrity relies on cryptographic immutability:
once a transaction is committed, unauthorized modification becomes
detectable. Strict immutability, however, can conflict with legitimate
requirements such as correcting erroneous records, removing sensitive
information for privacy or regulatory compliance, and updating obsolete
data~\cite{ref9,ref21,ref22,ref38}. The challenge is therefore not
simply how to modify blockchain data, but how to ensure that each
modification is \emph{policy-compliant, jointly authorized,
state-consistent, and auditable} without sacrificing scalability.

Consider a consortium blockchain receiving hundreds of compliance-driven
redaction requests across different transaction batches. Some requesters
may prove eligibility using private credentials. Even after validation,
processing requests independently repeats committee approvals, Merkle
updates, and chameleon-hash adaptations. Naive batching, however, risks
executing stale or unauthorized requests under authorization intended
for valid ones. Later, an auditor examining many redactions may again
verify identical authorization, transition, and checkpoint evidence.
Scalable redaction therefore requires safe sharing of cryptographic work
across \emph{validation, authorization, execution, and verification}
without expanding redaction authority.

Chameleon hashes are widely used to modify committed data while
preserving blockchain linkage~\cite{ref1,ref2,ref6}. Subsequent
schemes introduce fine-grained access control~\cite{ref8,ref15},
decentralized trapdoor management~\cite{ref3,ref19}, dynamic
updates~\cite{ref5,ref16}, threshold redaction~\cite{ref27}, and
traceable or publicly auditable redaction~\cite{ref33,ref34,ref35}.
Nevertheless, preserving a ledger commitment does not prove that a
modification was authorized for the particular requester, operation,
policy, and current state. Because an authorized chameleon-hash
collision intentionally preserves ledger consistency, conventional
chain verification alone cannot establish \emph{why} a redaction was
allowed or whether the authorized request matches the modified state.

Privacy and long-term security introduce further requirements.
Fine-grained and policy-hiding schemes restrict redaction authority or
conceal policies~\cite{ref4,ref18,ref31,ref34}, while verifiable designs
support redaction auditing~\cite{ref13,ref20,ref33}. However,
authorization may depend on private credentials that should not be
disclosed, and auditors should not require sensitive transaction
content. Meanwhile, quantum-resistant chameleon hashes have been
studied for redactable blockchains~\cite{ref10,ref11,ref17,ref37}.
Protecting only the chameleon hash is insufficient if requester
authentication, private-policy verification, committee approval, or
audit authentication remains quantum-vulnerable.

Scalability must therefore be addressed jointly with authorization and
verifiability. Prior work considers scalable redaction
~\cite{ref14,ref21,ref26}, efficient data-structure updates
~\cite{ref23}, lightweight verification~\cite{ref20}, and robust
threshold redaction~\cite{ref27}. Yet request-by-request processing
repeats policy, committee, Merkle, and chameleon-hash operations, while
conventional batching does not exploit whether requests affect the same
transaction batches or share Merkle paths. Similarly, independent
auditing repeatedly verifies evidence common to multiple redactions.
The key problem is therefore to \emph{identify, preserve, and safely
reuse cryptographic sharing} throughout the redaction lifecycle while
retaining request-level authorization and state binding.

To address this problem, we propose \emph{VeRedact-PQ}, a scalable
post-quantum and verifiable redaction framework for permissioned
blockchains. VeRedact-PQ first binds each request to its authenticated
transaction and policy state through PQ authentication and PQZK-based
private-policy validation. It then introduces a coordinated optimization
pipeline: \emph{Cost- and Overlap-Aware Redaction Aggregation (CORA)}
identifies validated requests that can share execution work;
\emph{Shared-Frontier Merkle Redaction (SFMR)} realizes this sharing
during state modification; and \emph{State-Aware Verification Evidence
Reuse (SAVER)} preserves and reuses shared dependencies during auditing.
A \emph{Policy-Bound Batch Redaction Proof (PBRP)} links validation,
committee authorization, and the finalized transition.

The main contributions are summarized as follows:

\begin{itemize}

\item \textbf{State-Bound Post-Quantum Redaction.}
We design an end-to-end security-critical redaction path integrating
PQ requester authentication, state-bound PQZK private-policy
verification, $t$-of-$n$ PQ committee authorization, and distributed
Post-Quantum Chameleon Hash (PQCH) adaptation. The
\emph{Redaction-Aware Dual Commitment (RADC)} separates immutable
transaction-policy context from mutable content state, while SA-RLI
supports authenticated state resolution.

\item \textbf{Cost- and Overlap-Aware Redaction Aggregation.}
We introduce \emph{CORA}, which groups already validated requests using
workload conditions and authenticated execution locality. Unlike
load-only batching, CORA considers affected transaction batches and
shared Merkle frontiers to reduce subsequent cryptographic work, while
bounded waiting preserves fairness without allowing locality to
determine authorization eligibility.

\item \textbf{Shared-Frontier Verifiable Redaction.}
We develop \emph{SFMR} to jointly execute authorized modifications
sharing Merkle authentication structure, processing overlapping paths
once and requiring one distributed PQCH adaptation per affected batch.
PBRP binds request-level validation, threshold authorization, the
authorization-bound SFMR transition, and finalized checkpoint,
preserving individual verifiability under coalesced execution.

\item \textbf{State-Aware Verification Evidence Reuse.}
We introduce \emph{SAVER}, which organizes redaction auditing as
dependency-aware verification over content-addressed request,
authorization, transition, and checkpoint evidence. Shared evidence is
verified once and reused only under an identical authenticated state
context. Challenged requests can selectively escalate to independent
PQZK verification without revalidating unaffected shared evidence,
while RAI provides authenticated historical query resolution.

\end{itemize}


    
\section{Related Work}
   \begin{table*}[t]
\centering
\caption{Comparison of VeRedact-PQ with Representative Redactable
Blockchain Schemes}
\label{tab:comparison}
\renewcommand{\arraystretch}{1.15}
\setlength{\tabcolsep}{4.0pt}
\begin{tabular}{lcccccc}
\hline
\textbf{Scheme} &
\textbf{PQ} &
\textbf{Threshold} &
\textbf{Policy-Bound} &
\textbf{Cost/Overlap-} &
\textbf{State-Bound} &
\textbf{Reusable} \\
&
\textbf{Security} &
\textbf{Authorization} &
\textbf{Authorization} &
\textbf{Aware Redaction} &
\textbf{Execution} &
\textbf{Auditing} \\
\hline

Scheme~\cite{ref1}
& $\times$ & $\checkmark$ & $\times$ & $\times$
& $\times$ & $\times$ \\

Scheme~\cite{ref5}
& $\times$ & $\triangle$ & $\checkmark$ & $\times$
& $\triangle$ & $\times$ \\

Scheme~\cite{ref13}
& $\times$ & $\times$ & $\times$ & $\times$
& $\triangle$ & $\triangle$ \\

Scheme~\cite{ref14}
& $\times$ & $\times$ & $\triangle$ & $\triangle$
& $\triangle$ & $\times$ \\

Scheme~\cite{ref15}
& $\times$ & $\checkmark$ & $\checkmark$ & $\times$
& $\triangle$ & $\times$ \\

Scheme~\cite{ref17}
& $\checkmark$ & $\times$ & $\triangle$ & $\times$
& $\times$ & $\times$ \\

Scheme~\cite{ref20}
& $\times$ & $\triangle$ & $\checkmark$ & $\triangle$
& $\triangle$ & $\triangle$ \\

Scheme~\cite{ref27}
& $\times$ & $\checkmark$ & $\times$ & $\times$
& $\triangle$ & $\times$ \\

Scheme~\cite{ref33}
& $\times$ & $\triangle$ & $\checkmark$ & $\triangle$
& $\triangle$ & $\triangle$ \\

Scheme~\cite{ref34}
& $\times$ & $\checkmark$ & $\checkmark$ & $\times$
& $\triangle$ & $\times$ \\

Scheme~\cite{ref39}
& $\checkmark$ & $\times$ & $\times$ & $\times$
& $\times$ & $\times$ \\

\hline
\textbf{VeRedact-PQ}
& $\checkmark$ & $\checkmark$ & $\checkmark$
& $\checkmark$ & $\checkmark$ & $\checkmark$ \\
\hline

\multicolumn{7}{l}{
\footnotesize $\checkmark$: Supported;
$\triangle$: Partially supported;
$\times$: Not supported. Scheme~\cite{ref39} is a
post-quantum batch authentication scheme without redaction.}
\end{tabular}
\end{table*}
    \subsection{Distributed and Scalable Blockchain Redaction}
    
    Redactable blockchains commonly employ chameleon hashes to modify
    committed data while preserving blockchain consistency~\cite{ref9}.
    Subsequent studies extended this model to consortium blockchains
    ~\cite{ref2} and decentralized redaction through distributed chameleon
    hashes and trapdoor management~\cite{ref1,ref3,ref6,ref19}.
    Multi-party chameleon hashing~\cite{ref12} and robust threshold
    redaction~\cite{ref27} further reduce dependence on a single redaction
    authority, while dynamic and controlled constructions support evolving
    redaction requirements~\cite{ref5,ref16,ref36}.
    
    Scalability has also been addressed through efficient update
    mechanisms, lightweight storage, specialized Merkle structures, and
    application-oriented redaction~\cite{ref14,ref20,ref21,ref23,ref26}.
    However, efficiency is typically optimized at a particular layer,
    whereas high-volume redaction incurs combined costs from target lookup,
    request validation, committee authorization, authenticated state
    updates, and chameleon-hash adaptation. VeRedact-PQ addresses these
    costs jointly through sharded lookup, adaptive request batching,
    coalesced Merkle updates, and one distributed PQCH adaptation per
    affected transaction-batch root transition.
    
    \subsection{Policy-Controlled, Private, and Verifiable Redaction}
    
    Fine-grained redaction has been developed using attribute- and
    policy-based mechanisms~\cite{ref8,ref15}, while privacy-preserving and
    policy-hiding constructions restrict disclosure of sensitive
    authorization information~\cite{ref4,ref18,ref24,ref31,ref34}. Other schemes
    strengthen traceability and accountability through dynamic policies,
    anonymous but accountable redaction, policy-compliant rewriting, and
    cross-chain accountability~\cite{ref25,ref28,ref29,ref32,ref35}. Verifiable
    redaction has likewise been studied through efficient query and
    integrity auditing~\cite{ref13}, lightweight verification and
    permission supervision~\cite{ref20}, and publicly auditable
    architectures~\cite{ref33}.
    
    Despite these advances, policy validation, multi-party authorization,
    redaction execution, and auditing are often treated as separate
    functions. VeRedact-PQ instead binds them at request level: private
    policy eligibility is verified before batching, committee authorization
    is bound to an authenticated batch, and PBRP links this evidence to the
    executed state transition. A sharded RAI subsequently supports
    query-based auditing using compact multiproofs and shared batch
    evidence, preserving individual accountability without repeatedly
    processing complete redaction proofs.
    
    \subsection{Post-Quantum Redactable Blockchains}
    
    Post-quantum redactable blockchain research has primarily focused on
    the chameleon-hash primitive. Existing approaches include
    quantum-resistant key-exposure-free chameleon hashing~\cite{ref11},
    lattice-based redactable blockchains~\cite{ref10}, quantum-resistant
    redaction with trapdoor updates~\cite{ref17}, and lattice-based policy-
    or identity-oriented chameleon hashes~\cite{ref18,ref37}. These works
    provide important foundations for protecting controlled redaction
    against quantum adversaries.
    
    However, a quantum-resistant chameleon hash alone does not provide an
    end-to-end post-quantum redaction workflow if authentication,
    authorization, committee approval, or auditing still relies on
    classical mechanisms. VeRedact-PQ therefore extends PQ protection
    across the security-critical workflow using PQ authentication,
    PQZK-based private policy validation, PQ multi-party authorization, and
    distributed PQCH redaction. To control overhead, expensive PQZK
    verification is performed only when required, while compact PQ-signed
    attestations and hash-based authenticated proofs support subsequent
    authorization and routine auditing.
    
    Overall, VeRedact-PQ differs from existing work by jointly addressing
    \emph{post-quantum security, scalable multi-party redaction, and
    privacy-preserving verifiability} across the complete redaction
    lifecycle rather than optimizing these properties independently
(Table~\ref{tab:comparison}).
    


 \section{Our Proposed VeRedact-PQ Scheme}

This section presents the system model, threat model, and protocol of
VeRedact-PQ.

\subsection{System Model}

\begin{figure}[t]
\centering
\includegraphics[width=1\linewidth]{VeRedact-PQ System Model.png}
\caption{System model of the proposed VeRedact-PQ framework}
\label{fig:system_model}
\end{figure}

As illustrated in Fig.~\ref{fig:system_model}, VeRedact-PQ considers a
permissioned blockchain operated by mutually accountable consortium
organizations. Transactions may subsequently require legitimate
modification or removal under predefined redaction policies. Rather
than assigning redaction capability to a single authority, VeRedact-PQ
separates request validation, threshold authorization, redaction
execution, and post-redaction auditing. The system comprises the
following entities.

\textbf{1) Data Owner (DO):}
The DO generates and signs transactions before submission to the
blockchain, providing origin authentication and integrity. Each
transaction is associated with an identifier, data type, timestamp, and
redaction policy. A DO may also act as an authorized requester when
permitted by the applicable policy.

\textbf{2) Authorized Requester (AR):}
An AR is an authenticated entity permitted to request modification or
removal of committed data. It submits a PQ-signed request together with
the evidence required to satisfy the applicable public and private
policy conditions. An AR may be the original DO or another entity
granted redaction rights by the consortium.

\textbf{3) Redaction Committee (RC):}
The RC, denoted $\mathcal{C}_e$, consists of $n$ consortium members
appointed for authorization epoch $e$. At least $t$ members must approve
a redaction batch. Committee members also hold distributed shares of
the PQCH trapdoor and participate in authorized PQCH adaptation.
Consequently, fewer than $t$ committee members cannot independently
exercise the redaction capability.

\textbf{4) Validation and Policy Service (VPS):}
The VPS authenticates requesters, resolves authenticated transaction
state, evaluates public policy conditions, and verifies state-bound
PQZK proofs for private policy conditions. Successfully validated
requests receive PQ-signed validation attestations and enter the
pending queue $\mathcal{Q}_e$. The VPS also issues signed admission
receipts that provide verifiable request waiting times for subsequent
CORA processing.

\textbf{5) Permissioned Blockchain Network (PBN):}
The PBN comprises authenticated consortium nodes executing the
underlying permissioned consensus protocol. It maintains transaction
and policy commitments, Merkle-batch states, PQCH commitments,
authenticated index roots, authorization evidence, and finalized
redaction checkpoints. Large objects, including PQZK proofs and PBRP
evidence, may be retained in an off-ledger evidence store, with their
digests or authenticated roots anchored on-chain.

The PBN additionally provides two service functions. A
\emph{checkpoint authority} signs finalized batch/index checkpoints,
while an \emph{audit service} maintains the Redaction Audit Index (RAI),
resolves authorized audit queries, and signs responses. Their outputs
remain independently verifiable against blockchain-anchored
commitments.

\textbf{6) Auditor (AU):}
An AU is an authorized internal or external verifier that examines
completed redactions without requiring access to the underlying
sensitive transaction content. It verifies request validation,
committee authorization, redaction-state transitions, PQCH commitment
preservation, and finalized blockchain evidence at the verification
depth required by the audit.

\subsection{Threat Model}
\label{subsec:threat}

VeRedact-PQ considers malicious or compromised participants in a
permissioned blockchain and quantum-capable adversaries. The main
threats and assumptions are as follows.

\textbf{1) Malicious Requester:}
A requester may submit unauthorized, malformed, conflicting, or replayed
redaction requests, modify data outside its policy scope, or reuse
evidence from another state. Requests are therefore PQ-authenticated
and validated against the applicable policy and authenticated
transaction state before CORA processing.

\textbf{2) Malicious Committee Members:}
Committee members may collude to authorize invalid redactions or misuse
PQCH shares. We assume fewer than $t$ members of the active committee
$\mathcal{C}_e$ are compromised. PQCH shares are proactively refreshed
or reshared across epochs over authenticated confidential channels,
with obsolete shares securely erased. Security against cross-epoch
accumulation assumes fewer than $t$ contemporaneously valid shares are
compromised within each resharing interval.

\textbf{3) Malicious Blockchain or Service Participant:}
A compromised consortium node or audit service may alter, omit, inject,
or replay transaction, checkpoint, index, or audit evidence. We assume
the PBN satisfies its prescribed consensus fault-tolerance condition.
Authenticated roots, PQ signatures, PQCH commitments, and finalized
checkpoints make inconsistent evidence detectable. Audit completeness
is claimed only when authenticated coverage evidence is available.

\textbf{4) Authorization and Execution Manipulation:}
An adversary may insert or substitute requests after CORA selection,
reuse authorization for another state, or execute unauthorized
positions. The validated-request commitments, $R_e^{VR}$, $\Gamma_e$,
and $\Theta_b$ bind authorization to its execution context. Phase~5
enforces
\begin{equation}
\mathcal{P}_b^{exec}\subseteq\mathcal{P}_b^{auth}
\label{eq:subset}
\end{equation}
allowing stale requests to be withheld but preventing unauthorized
positions from being introduced.

\textbf{5) Verification-Reuse Adversary:}
An adversary may attempt to reuse SAVER evidence under a different
authorization, transition, or checkpoint state. SAVER binds each
verified result to its content and authenticated state context through
$\chi(v)$; any change to the evidence or its bound dependencies
requires fresh verification.

\textbf{6) Privacy-Curious Auditor:}
An auditor may attempt to infer information beyond its authorized
scope. PQZK hides private credentials and attributes, while
query-scoped evidence limits unnecessary content disclosure. Access-
pattern privacy and confidentiality against PBN nodes storing the
transaction content are outside the scope of this work.

\textbf{7) Quantum-Capable Adversary:}
The security-critical redaction path uses PQ signatures, PQZK, PQCH,
and quantum-resistant hash/PRF assumptions. The underlying permissioned
consensus and platform credentials are not required to be
post-quantum.  \subsection{System Process}

VeRedact-PQ operates through six sequential phases covering system
initialization, transaction commitment, request validation, authorization,
redaction execution, and post-redaction auditing.

    
    \subsubsection*{\textbf{Phase 1: System Initialization}}

This phase initializes the cryptographic parameters, registered
entities, redaction policies, and distributed redaction authority of
VeRedact-PQ. Let $\lambda$ denote the security parameter and $e$ the
current authorization epoch.

\textbf{Step 1: Global Parameter Setup.}
Given $1^\lambda$, the consortium initializes the cryptographic
primitives and publishes
\begin{equation}
\begin{aligned}
PP=(&\lambda,H,H_1,H_2,H_A,H_V,\mathsf{PQCH},\\
&\mathsf{PQSIG},\mathsf{PQZK},\mathsf{Merkle},\mathsf{PRF})
\end{aligned}
\end{equation}
where $H$ is a post-quantum-secure hash function;
$H_1$, $H_2$, $H_A$, and $H_V$ are domain-separated instances of
$H$; $\mathsf{PQCH}$ is a post-quantum chameleon-hash scheme;
$\mathsf{PQSIG}$ is a post-quantum signature scheme;
$\mathsf{PQZK}$ is a post-quantum zero-knowledge proof system;
$\mathsf{Merkle}$ is an authenticated Merkle-tree construction; and
$\mathsf{PRF}$ is a quantum-secure pseudorandom function.
The additional domain $H_V$ is reserved for authenticated
verification-evidence commitments used by the state-aware verification
protocol in Phase~6.

\textbf{Step 2: Entity Registration and Key Generation.}
Each registered entity $U_i$, including data owners, authorized
requesters, committee members, the Validation and Policy Service
(VPS), PBN checkpoint and audit services, and auditors, generates
\begin{equation}
(pk_i,sk_i)\leftarrow
\mathsf{PQSIG.KeyGen}(1^\lambda)
\end{equation}
The tuple $(ID_i,Role_i,pk_i)$ is recorded in the consortium membership
registry, while $sk_i$ remains private. In particular, the VPS uses
$(pk_V,sk_V)$ for validation attestations; the checkpoint authority uses
$(pk_{\mathrm{CP}},sk_{\mathrm{CP}})$ for checkpoint signatures; and the
audit service uses $(pk_A,sk_A)$ for authenticated audit responses.
A committee member $C_k$ active in epoch $e$ uses its registered pair
$(pk_{e,k},sk_{e,k})$.

The consortium additionally samples
\begin{equation}
K_{\mathrm{idx}},K_A
\stackrel{\$}{\leftarrow}\{0,1\}^{\lambda}
\end{equation}
where $K_{\mathrm{idx}}$ derives SA-RLI lookup tokens and is available
to the PBN and VPS, whereas $K_A$ derives RAI audit tokens and is held
by the PBN audit service. The parameter $\lambda$ is selected according
to the target post-quantum security level.

\textbf{Step 3: Redaction Policy Registration.}
For each supported transaction or data class, the consortium defines a
redaction policy $P_j$ specifying eligible requester roles, permitted
operations, applicable conditions, and validity period. Its authenticated
state is committed as
\begin{equation}
C_{P_j}
=
H(PID_j\parallel P_j\parallel v_j\parallel e_P)
\label{eq:policy_commit}
\end{equation}
where $PID_j$, $v_j$, and $e_P$ denote the policy identifier, version,
and commitment epoch, respectively. The tuple
$(PID_j,C_{P_j},v_j,e_P)$ is registered on-chain and subsequently binds
request validation and redaction authorization to an authenticated
policy state.

\textbf{Step 4: Epoch Committee Formation.}
At the beginning of epoch $e$, the consortium establishes an
$n$-member redaction committee
\begin{equation}
\mathcal{C}_e=(C_1,C_2,\ldots,C_n)\qquad t\leq n
\end{equation}
where at least $t$ members are required to authorize and exercise the
distributed redaction capability. The configuration
$(e,\mathcal{C}_e,t)$ is recorded on-chain, binding subsequent
authorization and PQCH adaptation to the committee active in epoch $e$.

\textbf{Step 5: Distributed PQCH Trapdoor Setup.}
At system initialization, the initial committee
$\mathcal{C}_{e_0}$ jointly executes
\begin{equation}
(pk_{\mathrm{CH}},\{td_{e_0,k}\}_{k=1}^{n})
\leftarrow
\mathsf{PQCH.DKeyGen}(PP,\mathcal{C}_{e_0},t)
\end{equation}
where $pk_{\mathrm{CH}}$ is the public chameleon-hash key and
$td_{e_0,k}$ is the trapdoor share of member $C_k$.
The public key $pk_{\mathrm{CH}}$ is recorded on the PBN and remains
fixed across authorization epochs.

At each transition $e-1\rightarrow e$, the trapdoor is proactively
reshared to the incoming committee:
\begin{equation}
\{td_{e,k}\}_{k=1}^{n}
\leftarrow
\mathsf{PQCH.Reshare}
(\{td_{e-1,k}\},\mathcal{C}_e,t)
\end{equation}
The resharing protocol is executed over authenticated confidential
channels, and previous-epoch shares are securely erased after successful
resharing. Hence, the active committee can adapt existing batch
commitments without reconstructing the complete trapdoor at any single
member. Security against cross-epoch share accumulation assumes secure
erasure and that fewer than $t$ valid shares are compromised within
each resharing interval.

Threshold adaptation is specified as a functionality
$\mathcal{F}_{\mathrm{TA}}$: on input at least $t$ shares and an
authorized transition $(MR_b,r_b,MR_b')$, it outputs a valid collision
randomness $r_b'$ and nothing else. For an SIS-based PQCH with a gadget
trapdoor $R$~\cite{ref10,ref17}, where a collision is sampled as
$r_b'=p+[R;I]z$ for a perturbation $p$ and a gadget preimage $z$,
$\mathcal{F}_{\mathrm{TA}}$ is realized by evaluating the
trapdoor-dependent terms, $Rz$ and the perturbation, by secure
multi-party computation over the shares. No party, including the share
combiner, therefore learns $Rz$ or any other function of $R$ beyond the
output $r_b'$.

\textbf{Step 6: Post-Quantum ZK Parameter Setup.}
For the policy-compliance relation $\mathcal{R}_P$, the consortium
initializes
\begin{equation}
(pp_{\mathrm{PQZK}},pk_{\mathrm{PQZK}},vk_{\mathrm{PQZK}})
\leftarrow
\mathsf{PQZK.Setup}(1^\lambda,\mathcal{R}_P)
\end{equation}
The selected PQZK construction provides completeness, soundness, and
zero knowledge against quantum-capable adversaries. A transparent,
trapdoor-free setup is assumed; therefore, the public proving and
verification parameters can be published on the PBN without requiring
trusted setup-secret erasure or regeneration when new policy instances
are registered. Any prover-specific secret state required by the
selected instantiation remains with the corresponding prover.

All public keys, policy commitments, committee configurations, and
other security-critical initialization records are submitted as
transactions authenticated by the registering consortium members'
post-quantum signatures. Consequently, the cryptographic redaction
control plane does not rely solely on the underlying PBN's platform
credentials.
    
    
    
   \subsubsection*{\textbf{Phase 2: Redaction-Aware Transaction Commitment}}

This phase transforms each accepted transaction into a redaction-aware
authenticated state. VeRedact-PQ separates immutable transaction and
policy context from mutable content, commits transactions in
PQCH-protected Merkle batches, and constructs an authenticated lookup
index that exposes the batch and leaf locality required by subsequent
authorization and optimized redaction. Unless stated otherwise, the
operations in this phase are performed by the PBN.

\medskip
\noindent\textbf{Step 1: Transaction Authentication and Immutable
Provenance.}
A registered data owner $DO_i$ submits content $m_i$ with transaction
identifier $TID_i$, data type $DT_i$, redaction policy $PID_i$, and
timestamp $ts_i$. The owner samples
$\rho_i\stackrel{\$}{\leftarrow}\{0,1\}^{\lambda}$ and computes
\begin{equation}
D_i=H(m_i\parallel\rho_i)
\end{equation}
followed by
\begin{equation}
\begin{aligned}
\sigma_i\leftarrow
\mathsf{PQSIG.Sign}\big(
sk_{DO_i},
H(&TID_i\parallel D_i\parallel DT_i\\
&\parallel PID_i\parallel ts_i)\big)
\end{aligned}
\end{equation}
The submitted transaction is
\begin{equation}
TX_i=(TID_i,m_i,\rho_i,DT_i,PID_i,ts_i,\sigma_i)
\end{equation}
The PBN recomputes $D_i$ and verifies transaction uniqueness, the
registered owner key, $\sigma_i$, and the referenced policy before
acceptance. The verified key determines $ID_{DO_i}$.

To preserve the authenticated origin independently of later content
states, the PBN derives
\begin{equation}
C_i^{\mathrm{orig}}
=
H(TID_i\parallel D_i\parallel\sigma_i)
\end{equation}
and anchors the non-redactable provenance tuple
\begin{equation}
\mathcal{P}_i^{\mathrm{orig}}
=
(TID_i,DT_i,PID_i,ts_i,D_i,\sigma_i,C_i^{\mathrm{orig}})
\end{equation}
Hence, $\sigma_i$ continues to authenticate the originally submitted
content commitment after an authorized redaction; subsequent content
states are authenticated by the redaction evidence generated in
Phases~4--6 rather than by re-signing them as original owner content.

\medskip
\noindent\textbf{Step 2: Redaction-Aware Dual Commitment.}
For each accepted transaction, the PBN samples a fresh nonce $n_i$ and
computes
\begin{equation}
Tag_i=H(TID_i\parallel n_i\parallel ts_i)
\end{equation}
The Redaction-Aware Dual Commitment (RADC) separates immutable context
from mutable content through
\begin{equation}
\begin{aligned}
I_i=H(&TID_i\parallel ID_{DO_i}\parallel DT_i
\parallel PID_i\\
&\parallel ts_i\parallel Tag_i)
\end{aligned}
\end{equation}
and the content commitment $D_i$ from Step~1. The corresponding
authenticated leaf is
\begin{equation}
L_i=H(I_i\parallel D_i
\end{equation}

An authorized transition from $(m_i,\rho_i)$ to
$(m_i',\rho_i')$ therefore changes only
\begin{equation}
D_i\rightarrow D_i'=H(m_i'\parallel\rho_i')
\end{equation}
while $I_i$ remains invariant. This separation prevents a content
redaction from altering the transaction identity, ownership context,
data type, policy binding, or original provenance. The current content
and opening randomness are maintained in the transaction state, whereas
the immutable provenance record remains unchanged.

\medskip
\noindent\textbf{Step 3: Redaction-Locality Batch Commitment.}
Authenticated leaves are placed at deterministic positions in a
transaction batch
\begin{equation}
\mathcal{T}_b=(L_{b,1},\ldots,L_{b,N_b})
\end{equation}
whose root is
\begin{equation}
MR_b=
\mathsf{Merkle.Root}(L_{b,1},\ldots,L_{b,N_b})
\end{equation}
The batch is then committed (Fig.~\ref{fig:merkle}) using
\begin{equation}
CH_b=
\mathsf{PQCH.Hash}(pk_{\mathrm{CH}},MR_b,r_b)
\qquad
r_b\stackrel{\$}{\leftarrow}\mathcal{R}_{\mathrm{CH}}
\end{equation}
where $\mathcal{R}_{\mathrm{CH}}$ denotes the randomness space of the
selected PQCH scheme. The initial batch version is $v_b=0$.

The Merkle topology is retained as authenticated update state. For any
future set of leaves $\mathcal{L}_b^{*}$ selected for redaction, define
its minimal external authentication frontier as
\begin{equation}
\mathcal{F}_b(\mathcal{L}_b^{*})
=
\left\{
u:
u\notin\mathsf{Anc}(\mathcal{L}_b^{*}),
\;
\mathsf{parent}(u)\in
\mathsf{Anc}(\mathcal{L}_b^{*})
\right\}
\end{equation}
where $\mathsf{Anc}(\mathcal{L}_b^{*})$ is the union of Merkle nodes on
paths from the selected leaves to $MR_b$. The frontier is not computed
during initial commitment; rather, the stored topology permits it to be
derived when multiple authorized redactions are coalesced in Phase~5.
Consequently, overlapping authentication paths can later be represented
once instead of independently for every modified leaf.

\begin{figure}[t]
\centering
\includegraphics[width=0.7\linewidth]{Merkle.PNG}
\caption{Redaction-aware Merkle batch and PQCH commitment structure.}
\label{fig:merkle}
\end{figure}

\medskip
\noindent\textbf{Step 4: Authenticated Redaction-State Indexing.}
To resolve redaction targets without scanning blockchain transactions,
VeRedact-PQ constructs the Sharded Authenticated Redaction Lookup Index
(SA-RLI). For each transaction,
\begin{equation}
\tau_i=\mathsf{PRF}_{K_{\mathrm{idx}}}(TID_i)
\end{equation}
and the corresponding authenticated state entry is
\begin{equation}
E_i=
(\tau_i,Tag_i,b,pos_i,PID_i,e_i,I_i,D_i,ptr_i)
\end{equation}
where $b$ and $pos_i$ identify the transaction's batch and leaf
position, $e_i$ identifies the epoch of the indexed state, and $ptr_i$
references the authenticated transaction state required to resolve the
leaf. The batch version $v_b$ is obtained from the finalized checkpoint
rather than duplicated in every entry.

Each entry is deterministically assigned to a shard and logical bucket:
\begin{equation}
sid_i=H_1(\tau_i)\bmod S
\qquad
bid_i=H_2(\tau_i)\bmod B
\end{equation}
where $S$ and $B$ are the numbers of index shards and logical buckets.
Within each shard, entries are canonically ordered by
$(bid_i,\tau_i)$ and committed as
\begin{equation}
R_s=
\mathsf{Merkle.Root}(E_{s,1},\ldots,E_{s,n_s})
\end{equation}
with global authenticated index state
\begin{equation}
R_{\mathrm{RLI}}
=
\mathsf{Merkle.Root}(R_1,\ldots,R_S)
\end{equation}

The explicit $(b,pos_i)$ mapping is also used by the later
cost- and overlap-aware redaction protocol to identify requests that
share a transaction batch or overlapping Merkle paths without exposing
or modifying their authenticated content state.

\medskip
\noindent\textbf{Step 5: Atomic State Checkpointing.}
After the transaction batch and its SA-RLI updates are finalized, the
PBN assigns a monotonically increasing snapshot identifier
$v_{\mathrm{RLI}}$ and constructs
\begin{equation}
\begin{aligned}
A_b=\big(&b,CH_b,MR_b,r_b,R_{\mathrm{RLI}},
v_{\mathrm{RLI}},\\
& e,v_b,ts_b,\sigma_b\big)
\end{aligned}
\end{equation}
where $e$ is the checkpoint-finalization epoch. The checkpoint authority
computes
\begin{equation}
\begin{aligned}
\sigma_b=
\mathsf{PQSIG.Sign}\big(
sk_{\mathrm{CP}},
H(&b\parallel CH_b\parallel MR_b\parallel r_b\\
&\parallel R_{\mathrm{RLI}}\parallel v_{\mathrm{RLI}}
\parallel e\\
&\parallel v_b\parallel ts_b)\big)
\end{aligned}
\end{equation}
The signed checkpoint is anchored to the PBN, atomically binding the
PQCH-protected transaction-batch state to the corresponding finalized
SA-RLI snapshot. Thus, later authorization and redaction operate against
a uniquely identified tuple
\begin{equation}
\mathsf{State}_b=
(b,MR_b,CH_b,v_b,R_{\mathrm{RLI}},v_{\mathrm{RLI}},e)
\end{equation}
which prevents evidence from different batch or index versions from
being silently combined.

For lookup acceleration, each shard may maintain a Cuckoo filter
$CF_s$ over its current tokens. The filter is used only as a screening
structure: positive results require authenticated SA-RLI verification,
and a negative result is accepted only when $CF_s$ is synchronized with
the selected finalized snapshot; otherwise, authenticated lookup is
performed.    
    
   \subsubsection*{\textbf{Phase 3: Redaction Request Authentication and
PQZK-Based Policy Verification}}

This phase converts a redaction request into an authenticated,
policy-compliant, and state-bound validated request. VeRedact-PQ first
resolves the target against a finalized SA-RLI snapshot, authenticates
the requester, rejects requests that fail public policy checks, and
invokes PQZK only for the remaining private policy conditions. Accepted
requests receive a PQ-signed validation attestation and enter the pending
queue for cost- and overlap-aware authorization in Phase~4.

\medskip
\noindent\textbf{Step 1: Request Admission and Authenticated State
Resolution.}
A requester $U_r$ prepares
\begin{equation}
R_i=(ID_r,TID_i,op_i,D_i',e,ts_r,n_r)
\end{equation}
where $op_i$ is the requested redaction operation, $e$ is the
authorization epoch, and $(ts_r,n_r)$ provides freshness. The
corresponding redaction payload is
\begin{equation}
X_i^{\mathrm{red}}=(m_i',\rho_i')
\qquad
D_i'=H(m_i'\parallel\rho_i')
\end{equation}
so that $R_i$ commits to the proposed content state without exposing
$m_i'$ in the authorization record. The payload is retained for
execution in Phase~5, where its opening against $D_i'$ is verified
before any state transition.

After basic format, timestamp, nonce, and epoch screening, the VPS
derives $\tau_i$, $sid_i$, and $bid_i$ as in Phase~2, Step~4. A
synchronized Cuckoo filter may reject an absent token; otherwise, the
VPS performs authenticated SA-RLI lookup, retrieves $E_i$, verifies it
against the corresponding shard root and
$R_{\mathrm{RLI}}$, retrieves the signed checkpoint $A_b$, and verifies
\begin{equation}
L_i=H(I_i\parallel D_i)
\end{equation}
at position $pos_i$ against $MR_b$. It also verifies the checkpoint
signature and consistency between the batch state and finalized index
snapshot. The resolved state is $\mathsf{State}_b$ of Phase~2, Step~5,
where $e_b$ denotes the checkpoint-finalization epoch, distinguished
from the current authorization epoch $e$. The indexed epoch $e_i$
remains transaction-state metadata and need not equal either $e_b$ or
$e$.

The pair $(b,pos_i)$ is retained as authenticated locality metadata.
It is used only after validation to estimate common-batch and
Merkle-path overlap in Phase~4 and does not affect request eligibility
or policy satisfaction.

\medskip
\noindent\textbf{Step 2: Post-Quantum Requester Authentication.}
The requester signs the complete request:
\begin{equation}
\sigma_i^R\leftarrow
\mathsf{PQSIG.Sign}(sk_r,H(R_i))
\end{equation}
and submits $(R_i,\sigma_i^R)$. Using the registered key $pk_r$, the
VPS verifies
\begin{equation}
\mathsf{PQSIG.Verify}
(pk_r,H(R_i),\sigma_i^R)=1
\label{eq:req_sig}
\end{equation}
together with the binding between $ID_r$ and $pk_r$, requester
registration, and eligibility to submit requests in epoch $e$.
The nonce $n_r$ is consumed in the authenticated requester's replay
domain only after successful authentication.

Only after authentication does the VPS release to the requester the
authenticated policy and state information required to construct the
PQZK statement.

\medskip
\noindent\textbf{Step 3: Public Policy and State Validation.}
Using the authenticated $PID_i$, the VPS retrieves policy $P_i$, its
version $v_i$, commitment epoch $e_P$, and on-chain commitment
$C_{P_i}$, and verifies
\begin{equation}
C_{P_i}
\stackrel{?}{=}
H(PID_i\parallel P_i\parallel v_i\parallel e_P)
\end{equation}
The VPS applies the registered policy-version and epoch-transition rules
to confirm that $P_i$ governs the resolved transaction state and is
valid for authorization epoch $e$. It then verifies $op_i$, all public
policy predicates, request freshness, and the validity of the resolved
state $\mathsf{State}_b$.

Requests failing these checks are rejected before PQZK verification.
This staged validation avoids invoking the comparatively expensive
PQZK verifier for requests that already fail authenticated public
conditions.

\medskip
\noindent\textbf{Step 4: State-Bound PQZK Policy Verification.}
For requests that pass the public checks, the requester proves
satisfaction of the private policy predicates without revealing its
credentials or sensitive authorization attributes. To bind the proof
to the exact authenticated state, define
\begin{equation}
\begin{aligned}
h_i^{\mathrm{leaf}} &= H\big(TID_i\parallel I_i\parallel D_i\parallel pos_i\big)\\
h_i^{\mathrm{snap}} &= H\big(\mathsf{enc}(\mathsf{State}_b)\big)
\end{aligned}
\label{eq:hleaf_hsnap}
\end{equation}
\begin{equation}
h_i^{\mathrm{st}}
=
H\big(h_i^{\mathrm{leaf}}\parallel h_i^{\mathrm{snap}}\big)
\end{equation}
where $\mathsf{enc}(\cdot)$ denotes an unambiguous canonical encoding.
The leaf part $h_i^{\mathrm{leaf}}$ fixes what the request
is about (transaction, immutable context, current content commitment and
position); the snapshot part $h_i^{\mathrm{snap}}$ records the finalized
checkpoint of batch $b$ against which the request is validated. The
requester submits $(h_i^{\mathrm{leaf}},h_i^{\mathrm{snap}})$ together
with $R_i$.
The public statement is
\begin{equation}
\begin{aligned}
x_i=\big(&ID_r,TID_i,op_i,D_i',
C_{P_i},h_i^{\mathrm{st}},\\
& e_i,e,ts_r,n_r\big)
\end{aligned}
\label{eq:pqzk_stmt}
\end{equation}
and the witness $w_i$ contains the private credentials,
authorization attributes, and policy evidence.

The policy relation is
\begin{equation}
\begin{aligned}
\mathcal{R}_{P}(x_i,w_i)=1
\iff{}&
\mathsf{ValidCred}(w_i,C_{P_i},e)\\
&\land
\mathsf{RequesterBound}(w_i,ID_r)\\
&\land
\mathsf{PrivatePolicy}(x_i,w_i)
\end{aligned}
\end{equation}
where credential validity includes the applicable expiry and revocation
conditions, and requester binding proves credential ownership or valid
delegation.

The requester computes
\begin{equation}
\pi_i^{PQ}\leftarrow
\mathsf{PQZK.Prove}
(pk_{\mathrm{PQZK}},x_i,w_i)
\end{equation}
and the VPS accepts only if
\begin{equation}
\mathsf{PQZK.Verify}
(vk_{\mathrm{PQZK}},x_i,\pi_i^{PQ})=1
\label{eq:pqzk_verify}
\end{equation}
The state digest $h_i^{\mathrm{st}}$ prevents a valid proof from being
transplanted to a different transaction position, content state, batch
version, or finalized SA-RLI snapshot.

Freshness is decided per target leaf rather than per
batch. Let $v_b^{cur}$ denote the current version of batch $b$ and $k$ a
system-wide snapshot window. A request is fresh if
\begin{equation}
\begin{aligned}
\mathsf{FreshLeaf}_i=1
\iff{}&
H(h_i^{\mathrm{leaf}}\parallel h_i^{\mathrm{snap}})=h_i^{\mathrm{st}}\\
&\land\
h_i^{\mathrm{leaf}}=H(TID_i\parallel I_i\parallel D_i^{cur}\parallel pos_i)\\
&\land\
h_i^{\mathrm{snap}}=H\big(\mathsf{enc}(\mathsf{State}_b^{(v)})\big)\\
&\phantom{\land\ }
v_b^{cur}-k\le v\le v_b^{cur}
\end{aligned}
\label{eq:freshleaf}
\end{equation}
where $\mathsf{State}_b^{(v)}$ is the state of the
anchored checkpoint $A_b$ of version $v$. The VPS accepts the request only
if $\mathsf{FreshLeaf}_i=1$, in addition to the epoch and policy-version
checks of Step~3. A redaction of another leaf of batch $b$ changes
$\mathsf{State}_b$ but not $h_i^{\mathrm{leaf}}$, and therefore no longer
invalidates pending requests on the same batch; a redaction of the target
leaf itself changes $D_i^{cur}$ and still does. Hence a fresh request is
executed only while its target leaf holds the content commitment against
which it was validated.

\medskip
\noindent\textbf{Step 5: Validated-Request Attestation and Queue
Admission.}
For each accepted request, the VPS assigns a unique $RID_i$ and derives
the proof digest
\begin{equation}
h_i^{PQ}=H(\pi_i^{PQ})
\end{equation}
It then constructs the validated-request commitment
\begin{equation}
\begin{aligned}
C_i^{VR}
=
H\big(&RID_i\parallel H(R_i)
\parallel h_i^{\mathrm{st}}
\parallel C_{P_i}\\
&\parallel h_i^{PQ}
\parallel e\big)
\end{aligned}
\label{eq:cvr}
\end{equation}
and admits it to the current admission window $w$ at time
$ts_i^{\mathrm{adm}}$ with the leaf
\begin{equation}
\nu_i=
H_V\big(\texttt{VR-LEAF}\parallel RID_i\parallel C_i^{VR}
\parallel ts_i^{\mathrm{adm}}\parallel e\big)
\end{equation}
A window closes after $W$ admissions or $T_{\mathrm{vps}}$
time units, whichever comes first. The VPS then commits its leaves as
$R_w^{V}=\mathsf{Merkle.Root}(\lambda_1,\ldots,\lambda_{|w|})$ and
signs once per window:
\begin{equation}
\sigma_w^{V}=
\mathsf{PQSIG.Sign}
\left(
sk_V,
H_V(\texttt{VPS-WIN}\parallel w\parallel R_w^{V}\parallel e)
\right)
\end{equation}
The attestation of request $i$ is
$\alpha_i=(w,R_w^{V},MP_i^{V},\sigma_w^{V},ts_i^{\mathrm{adm}})$,
where $MP_i^{V}$ is the membership proof of $\nu_i$ under
$R_w^{V}$. It is accepted iff
\begin{equation}
\begin{aligned}
&\mathsf{AttVerify}(\alpha_i)=1 \iff\\
&\quad \mathsf{PQSIG.Verify}\big(pk_V,\\
&\qquad H_V(\texttt{VPS-WIN}\parallel w\parallel R_w^{V}\parallel e),\,
\sigma_w^{V}\big)=1\\
&\quad \land\ \mathsf{Merkle.Verify}(\nu_i,MP_i^{V},R_w^{V})=1
\end{aligned}
\label{eq:attverify}
\end{equation}
A verifier checks $\sigma_w^{V}$ once per distinct
$(w,R_w^{V})$ and $\lceil\log_2 W\rceil$ hashes per request, instead of
one PQ signature per request.

The compact validated request is
\begin{equation}
\begin{aligned}
VR_i=\big(&RID_i,C_i^{VR},b,pos_i,PID_i,
v_b,v_{\mathrm{RLI}},\\
&e_i,e,h_i^{\mathrm{st}},h_i^{PQ},
\sigma_i^R,\alpha_i\big)
\end{aligned}
\end{equation}
The explicit $(b,pos_i)$ fields provide the authenticated locality
metadata required by Phase~4, whereas
$(v_b,v_{\mathrm{RLI}},h_i^{\mathrm{st}})$ bind that metadata to the
validated state. The complete request $R_i$, statement $x_i$, proof
$\pi_i^{PQ}$, policy evidence, and authenticated transaction evidence
remain available for independent or challenged verification.

The validated request is admitted to the pending queue
$\mathcal{Q}_e$ at time $ts_i^{\mathrm{adm}}$. The VPS returns
the receipt $rc_i=\alpha_i$: because $\nu_i$ binds
$ts_i^{\mathrm{adm}}$, the window signature and membership proof
provide verifiable evidence of admission. A request that reaches
the maximum waiting bound $T_{\max}$ becomes mandatory for the next
eligible authorization round, subject to state and policy revalidation.    
    
  \subsubsection*{\textbf{Phase 4: Cost- and Overlap-Aware Redaction
Aggregation and Multi-Party Authorization}}

This phase applies \emph{Cost- and Overlap-Aware Redaction Aggregation
(CORA)} to validated requests. Unlike load-only batching, CORA exploits
authenticated batch and leaf locality to co-schedule requests that can
share Merkle-update work and PQCH adaptation, while enforcing a bounded
waiting time. The active committee then revalidates the selected state
and issues $t$-of-$n$ post-quantum authorization over the resulting
batch commitment.

\medskip
\noindent\textbf{Step 1: CORA Candidate Selection.}
Let $\mathcal{Q}_e$ denote the queue of Phase~3 validated requests for
authorization epoch $e$. For each $VR_i\in\mathcal{Q}_e$, the authenticated
pair $(b_i,pos_i)$ identifies its target transaction batch and Merkle
position. For a candidate set $\mathcal{S}\subseteq\mathcal{Q}_e$, define
\begin{equation}
\Omega(\mathcal{S})
=
\{b_i:VR_i\in\mathcal{S}\}
\end{equation}
and, for each $b\in\Omega(\mathcal{S})$,
\begin{equation}
\mathcal{L}_b(\mathcal{S})
=
\{pos_i:VR_i\in\mathcal{S},\,b_i=b\}
\end{equation}
Using the fixed Merkle topology from Phase~2, the VPS derives the
corresponding external frontier
$\mathcal{F}_b(\mathcal{L}_b(\mathcal{S}))$. CORA estimates the
redaction cost as
\begin{equation}
\mathcal{C}_{\mathrm{CORA}}(\mathcal{S})
=
w_{\mathrm{CH}}|\Omega(\mathcal{S})|
+
w_{\mathrm{F}}
\sum_{b\in\Omega(\mathcal{S})}
|\mathcal{F}_b(\mathcal{L}_b(\mathcal{S}))|
\end{equation}
where $w_{\mathrm{CH}}$ and $w_{\mathrm{F}}$ reflect the relative cost
of a distributed PQCH adaptation and an authenticated frontier node,
respectively.

Given a workload-dependent target size
\begin{equation}
\widehat{B}_e
=
\min\!\left\{
B_{\max},
\max\!\left\{
B_{\min},
\left\lceil
f(\Lambda_e,|\mathcal{Q}_e|)
\right\rceil
\right\}
\right\}
\end{equation}
CORA selects
\begin{equation}
\mathcal{B}_e^{\mathrm{cand}}
=
\arg\min_{\substack{
\mathcal{S}\subseteq\mathcal{Q}_e\\
|\mathcal{S}|\leq\widehat{B}_e\\
\mathcal{M}_e\subseteq\mathcal{S}
}}
\mathcal{C}_{\mathrm{CORA}}(\mathcal{S})
\end{equation}
where
\begin{equation}
\mathcal{M}_e=
\{VR_i\in\mathcal{Q}_e:
t_{\mathrm{now}}-ts_i^{\mathrm{adm}}\geq T_{\max}\}
\end{equation}
is the mandatory set of requests reaching the waiting bound. If
$|\mathcal{M}_e|>\widehat{B}_e$, the oldest mandatory requests are
processed first and the target expands up to $B_{\max}$ when capacity
permits. Thus, locality can improve execution efficiency but cannot
indefinitely postpone an admitted request.

The combinatorial objective need not be solved exhaustively. The
implementation uses a deterministic greedy realization that admits
mandatory requests first and then repeatedly selects the request with
the smallest marginal cost
\begin{equation}
\Delta_i(\mathcal{S})
=
\mathcal{C}_{\mathrm{CORA}}
(\mathcal{S}\cup\{VR_i\})
-
\mathcal{C}_{\mathrm{CORA}}(\mathcal{S})
\end{equation}
with admission time and $RID_i$ used as deterministic tie-breakers.

\medskip
\noindent\textbf{Step 2: Attestation and Current-State Revalidation.}
Before authorization, each selected request is independently
revalidated. The committee verifies
$\mathsf{AttVerify}(\alpha_i)=1$
(Eq.~\eqref{eq:attverify}), verifying each distinct window signature
$\sigma_w^{V}$ once for the selected set,
and reconstructs $C_i^{VR}$ from the retained Phase~3 evidence.
Routine authorization verifies the VPS attestation , which
binds the proof digest $h_i^{PQ}$ through $C_i^{VR}$,
rather than repeating the underlying PQZK verification.

The committee then resolves the latest authenticated state
$\mathsf{State}_b^{cur}$ and re-evaluates
Eq.~\eqref{eq:freshleaf} against it.
The request remains eligible only if
$\mathsf{FreshLeaf}_i=1$,
epoch $e$ remains active, and the attested policy version remains
applicable. Requests that fail revalidation are removed and returned
for Phase~3 revalidation. CORA is not rerun merely to replace removed
requests; this prevents post-selection state changes from altering the
meaning of the candidate authorization round.

Let the remaining eligible requests, canonically ordered by
$(b_i,pos_i,RID_i)$, be
\begin{equation}
\mathcal{B}_e=(VR_1,\ldots,VR_m)
\end{equation}

\medskip
\noindent\textbf{Step 3: State-Bound Authorization Batch Commitment.}
For each $VR_i\in\mathcal{B}_e$, define
\begin{equation}
\eta_i
=
H_V\big(
\texttt{AUTH-LEAF}\parallel
RID_i\parallel C_i^{VR}\parallel
\alpha_i
\big)
\label{eq:eta}
\end{equation}
where the domain-separated $H_V$ from Phase~1 prevents
cross-context reuse. The leaves are committed as
\begin{equation}
R_e^{VR}
=
\mathsf{Merkle.Root}
(\eta_1,\ldots,\eta_m)
\end{equation}
Because $C_i^{VR}$ already binds the request, policy commitment,
PQZK-proof digest, authorization epoch, and authenticated state digest,
$\eta_i$ need not duplicate these fields.

To bind authorization to the execution structure selected by CORA,
define
\begin{equation}
\Gamma_e=
H_V\!\left(
\texttt{CORA-PLAN}\parallel
\mathsf{enc}
\left(
\{(b_i,pos_i)\}_{i=1}^{m}
\right)
\right)
\label{eq:cora_plan}
\end{equation}
and construct
\begin{equation}
C_e^{B}
=
H_V\big(
\texttt{AUTH-BATCH}\parallel
BID_e\parallel e\parallel m\parallel
R_e^{VR}\parallel\Gamma_e\parallel ts_e
\big)
\label{eq:auth_batch}
\end{equation}
where $BID_e$ uniquely identifies the authorization batch.
The plan commitment $\Gamma_e$ fixes the authorized batch/leaf
locality without committing to a particular execution implementation.

\medskip
\noindent\textbf{Step 4: Multi-Party Committee Authorization.}
Each approving member $C_k\in\mathcal{C}_e$ independently verifies the
finalized request set, current-state checks, $R_e^{VR}$, and
$\Gamma_e$, and signs
\begin{equation}
\sigma_{e,k}^{B}
\leftarrow
\mathsf{PQSIG.Sign}
\left(
sk_{e,k},C_e^{B}
\right)
\end{equation}
The valid approval set is
\begin{equation}
\begin{aligned}
\mathcal{A}_e=\big\{&
(k,\sigma_{e,k}^{B}):
C_k\in\mathcal{C}_e,\\
&
\mathsf{PQSIG.Verify}
(pk_{e,k},C_e^{B},\sigma_{e,k}^{B})=1
\big\}
\end{aligned}
\end{equation}
Authorization succeeds only when approvals from at least $t$ distinct
registered committee members are present:
\begin{equation}
|\mathcal{A}_e|\geq t
\end{equation}
The resulting shared authorization evidence is
\begin{equation}
Auth_e^{B}
=
(BID_e,C_e^{B},R_e^{VR},\Gamma_e,
\mathcal{A}_e,e,ts_e)
\end{equation}
This construction uses independent post-quantum signatures and does not
require a native threshold-signature scheme.

\medskip
\noindent\textbf{Step 5: Policy-Bound Authorization Evidence.}
For each authorized request, the system derives a Merkle membership
proof $MP_i^{B}$ for $\eta_i$ under $R_e^{VR}$ and constructs
\begin{equation}
PBRP_i^{auth}
=
(RID_i,\eta_i,MP_i^{B},Auth_e^{B})
\end{equation}
Since $Auth_e^{B}$ already contains $C_e^{B}$ and $R_e^{VR}$, they are
not duplicated in each authorization component. Shared committee
evidence can therefore be verified once per authorization batch, while
$MP_i^{B}$ binds each request to that evidence.    
 \subsubsection*{\textbf{Phase 5: Shared-Frontier Redaction and
Verifiable State Transition}}

This phase executes the CORA-authorized requests and converts their
shared locality into cryptographic execution savings. For each affected
transaction batch, VeRedact-PQ applies \emph{Shared-Frontier Merkle
Redaction (SFMR)} to recompute common Merkle paths once, performs one
threshold PQCH adaptation for the resulting batch-root transition, and
atomically binds the new transaction and SA-RLI states to the Phase~4
authorization.

\medskip
\noindent\textbf{Step 1: Authorization-Bound Execution Validation.}
Given the authorized request set $\mathcal{B}_e$, the PBN verifies
$Auth_e^B$, the threshold approval set $\mathcal{A}_e$, and each
$PBRP_i^{auth}$. It reconstructs the authorized position set and verifies
\begin{equation}
\Gamma_e
\stackrel{?}{=}
H_V\!\left(
\texttt{CORA-PLAN}\parallel
\mathsf{enc}\big(\{(b_i,pos_i)\}_{VR_i\in\mathcal{B}_e}\big)
\right)
\end{equation}
The requests are then partitioned by target batch:
\begin{equation}
\mathcal{B}_e
=
\bigcup_{b\in\Omega_e}\mathcal{B}_{e,b}
\qquad
\mathcal{B}_{e,b}
=
\{VR_i\in\mathcal{B}_e:b_i=b\}
\end{equation}
where $\Omega_e$ is the set of affected transaction batches.

Immediately before execution, the PBN resolves the latest authenticated
state for every request and re-evaluates
Eq.~\eqref{eq:freshleaf}.
Execution requires
$\mathsf{FreshLeaf}_i=1$,
together with a still-valid authorization epoch, policy, and approved
operation. Execution always starts from the latest
finalized state of batch $b$, i.e., the current $MR_b$, $r_b$ and $v_b$
used in Steps~3--5, not from the snapshot bound by $h_i^{\mathrm{snap}}$;
the snapshot only certifies that validation used an authentic, recent
checkpoint. Requests targeting the same transaction are conflict-checked;
at most one resulting content state is executed for a transaction in
one authorization round. Any stale, conflicting, or otherwise invalid
request causes the corresponding transition to be withheld rather than
silently replaced by an unauthorized request.

\medskip
\noindent\textbf{Step 2: RADC Transition Construction.}
For each executable $VR_i$, the PBN retrieves
\begin{equation}
X_i^{\mathrm{red}}=(m_i',\rho_i')
\end{equation}
and verifies
\begin{equation}
H(m_i'\parallel\rho_i')
\stackrel{?}{=}D_i'
\end{equation}
A valid request induces the RADC transition
\begin{equation}
(I_i,D_i)\longrightarrow(I_i,D_i')
\end{equation}
with
\begin{equation}
L_i=H(I_i\parallel D_i)
\qquad
L_i'=H(I_i\parallel D_i')
\end{equation}
Thus, execution modifies only the authorized content commitment while
preserving the immutable transaction-policy binding $I_i$ and the
original provenance established in Phase~2.

For batch $b$, let
\begin{equation}
\mathcal{P}_b=
\{pos_i:VR_i\in\mathcal{B}_{e,b}\}
\end{equation}
denote the authorized executable positions, and let
$\mathcal{L}_b$ and $\mathcal{L}_b'$ denote their old and replacement
leaves, respectively.

\medskip
\noindent\textbf{Step 3: Shared-Frontier Merkle Redaction.}
Rather than verifying and updating one authentication path per
redaction, SFMR derives the minimal external frontier introduced in
Phase~2:
\begin{equation}
\mathcal{F}_b
=
\mathcal{F}_b(\mathcal{P}_b)
\end{equation}
The frontier contains only unaffected sibling nodes required to
reconstruct the paths from all modified leaves to $MR_b$. The PBN first
verifies
\begin{equation}
\mathsf{SFMR.Verify}
(MR_b,\mathcal{P}_b,\mathcal{L}_b,\mathcal{F}_b)=1
\end{equation}
and then computes
\begin{equation}
MR_b'
=
\mathsf{SFMR.Update}
(\mathcal{P}_b,\mathcal{L}_b',
\mathcal{F}_b)
\end{equation}
Internal nodes shared by multiple modified paths are recomputed once.
Hence, the authentication material is proportional to the unique
frontier induced by the selected positions rather than to the sum of
independent Merkle paths. In particular, for $q_b=|\mathcal{P}_b|$,
the independent-path upper cost $O(q_b\log N_b)$ is replaced by work
proportional to the size of the induced authentication subtree and its
frontier, with the largest savings when CORA selects structurally
overlapping positions.

To bind the resulting root to the authorized execution, define
\begin{equation}
\begin{aligned}
\Theta_b
=
H_V\big(&\texttt{SFMR-TRANS}\parallel BID_e\parallel b
\parallel v_b\\
&\parallel MR_b\parallel MR_b'
\parallel \mathsf{enc}(\mathcal{P}_b)
\parallel \Gamma_e\big)
\end{aligned}
\label{eq:theta}
\end{equation}

\medskip
\noindent\textbf{Step 4: Authorization-Gated PQCH Adaptation.}
Before releasing a PQCH adaptation share, each participating committee
member $C_k$ independently verifies $Auth_e^B$, $\Gamma_e$, the
authorized membership of every request affecting batch $b$, and the
SFMR transition committed by $\Theta_b$. A share is generated only if
$\mathcal{P}_b$ is a subset of the positions authorized by
$\Gamma_e$ and $MR_b'$ is the deterministic SFMR result for the
corresponding authorized RADC transitions.

Member $C_k$ then contributes
\begin{equation}
\delta_{b,k}
\leftarrow
\mathsf{PQCH.PartAdapt}
(td_{e,k},MR_b,r_b,MR_b',\Theta_b)
\end{equation}
and at least $t$ validated contributions are combined into
\begin{equation}
r_b'
\leftarrow
\mathsf{PQCH.Combine}
(\{\delta_{b,k}\}_{k\in\mathcal{T}_b})
\qquad
|\mathcal{T}_b|\geq t
\end{equation}
Both steps run inside $\mathcal{F}_{\mathrm{TA}}$ (Phase~1, Step~5), so
the combination reveals only $r_b'$. The adaptation is accepted only if
\begin{equation}
\begin{aligned}
\mathsf{PQCH.Hash}(pk_{\mathrm{CH}},MR_b,r_b)
&=
\mathsf{PQCH.Hash}(pk_{\mathrm{CH}},MR_b',r_b')\\
&=CH_b
\end{aligned}
\label{eq:pqch_cont}
\end{equation}

Thus, all authorized modifications within batch $b$ produce one SFMR
root transition and one distributed PQCH adaptation. Distinct affected
batches can be prepared in parallel.

\medskip
\noindent\textbf{Step 5: Atomic Authenticated State Finalization.}
For every successfully prepared batch transition,
\begin{equation}
v_b'=v_b+1
\end{equation}
Each executed transaction updates its SA-RLI entry to
\begin{equation}
E_i'
=
(\tau_i,Tag_i,b,pos_i,PID_i,e,I_i,D_i',ptr_i')
\end{equation}
where $ptr_i'$ resolves the current authenticated state under $MR_b'$.
Unmodified transaction states remain unchanged.

Let $\Omega_S$ denote the SA-RLI shards containing modified entries.
Only these shards are recomputed:
\begin{equation}
R_s^{*}=
\begin{cases}
R_s', & s\in\Omega_S\\
R_s,  & s\notin\Omega_S
\end{cases}
\end{equation}
yielding
\begin{equation}
R_{\mathrm{RLI}}'
=
\mathsf{Merkle.Root}(R_1^{*},\ldots,R_S^{*})
\qquad
v_{\mathrm{RLI}}'=v_{\mathrm{RLI}}+1
\end{equation}

For each affected batch, the checkpoint authority constructs
\begin{equation}
\begin{aligned}
A_b'=\big(&b,CH_b,MR_b',r_b',
R_{\mathrm{RLI}}',v_{\mathrm{RLI}}'\\
&e,v_b',ts_b'\big)
\end{aligned}
\end{equation}
with checkpoint digest
\begin{equation}
\begin{aligned}
h_b^{\mathrm{CP}}=
H_V\big(&\texttt{CP}\parallel b\parallel CH_b\parallel MR_b'
\parallel H(r_b')\\
&\parallel R_{\mathrm{RLI}}'\parallel v_{\mathrm{RLI}}'
\parallel e\parallel v_b'\parallel ts_b'\big)
\end{aligned}
\end{equation}
Because the PBN finalizes the affected batches of a round
atomically, the checkpoint authority signs the round rather than each
batch: it commits the digests as
$R_e^{\mathrm{CP}}=\mathsf{Merkle.Root}(\{h_b^{\mathrm{CP}}\}_{b\in\Omega_e})$
and computes
\begin{equation}
\begin{aligned}
\sigma_e^{\mathrm{CP}}={}&
\mathsf{PQSIG.Sign}\big(sk_{\mathrm{CP}},\\
&\quad H_V(\texttt{CP-ROUND}\parallel BID_e\parallel R_e^{\mathrm{CP}}
\parallel v_{\mathrm{RLI}}')\big)
\end{aligned}
\end{equation}
Let $\Pi_b^{\mathrm{CP}}$ denote the membership proof of
$h_b^{\mathrm{CP}}$ under $R_e^{\mathrm{CP}}$. The finalization
transaction carries the $|\Omega_e|$ checkpoints, $R_e^{\mathrm{CP}}$, and
the single signature $\sigma_e^{\mathrm{CP}}$; validators verify
$\sigma_e^{\mathrm{CP}}$ once and recompute $R_e^{\mathrm{CP}}$ before
anchoring, and the PBN stores $h_b^{\mathrm{CP}}$ for each $(b,v_b')$. The PQCH randomness $r_b'$ enters the digest through $H(r_b')$; it remains
bound to $CH_b$ and $MR_b'$ by the PQCH equation.
All affected checkpoints reference the same finalized SA-RLI snapshot.
The PBN atomically commits the prepared batch and index transitions
only if their base states remain current; otherwise, the prepared
transitions are discarded and no new finalized state is published.

\medskip
\noindent\textbf{Step 6: Transition-Bound PBRP Generation.}
For each successfully executed request, define the shared batch
execution commitment
\begin{equation}
\begin{aligned}
\Xi_b=
H_V\big(&\texttt{EXEC-EVIDENCE}\parallel
\Theta_b\parallel CH_b\\
&\parallel v_b'\parallel
R_{\mathrm{RLI}}'\parallel v_{\mathrm{RLI}}'\big)
\end{aligned}
\label{eq:xi}
\end{equation}
The final Policy-Bound Batch Redaction Proof is
\begin{equation}
\begin{aligned}
PBRP_i=\big(
&PBRP_i^{auth},I_i,D_i,D_i',
b,pos_i,v_b,v_b',v_i^{\mathrm{snap}},\\
&\Theta_b,\Xi_b,A_b',\Pi_b^{\mathrm{CP}},\Pi_i^{\mathrm{SF}}\big)
\end{aligned}
\label{eq:pbrp}
\end{equation}
where $\Pi_i^{\mathrm{SF}}$ is the request-specific membership evidence
linking $(pos_i,L_i,L_i')$ to the shared SFMR transition,
and $v_i^{\mathrm{snap}}$ is the version of the checkpoint bound by
$h_i^{\mathrm{snap}}$, with $v_b-k\le v_i^{\mathrm{snap}}\le v_b$; here
$v_b$ is the version at execution, which may be newer than
$v_i^{\mathrm{snap}}$. The frontier,
PQCH adaptation evidence, and finalized checkpoint are batch-shared
objects and need not be duplicated across PBRPs.

The corresponding compact redaction record is
\begin{equation}
\begin{aligned}
RR_i=\big(
&RID_i,TID_i,BID_e,PID_i,D_i,D_i',\\
&v_b,v_b',H_V(\texttt{PBRP}\parallel
\mathsf{enc}(PBRP_i)),ts_i^{red}\big)
\end{aligned}
\end{equation}
Only successfully finalized transitions produce $RR_i$. The complete
PBRP and its referenced shared evidence are retained for Phase~6
verification and auditing.


\subsubsection*{\textbf{Phase 6: SAVER-Based State-Aware
Verifiable Redaction Auditing}}

This phase introduces \emph{State-Aware Verification Evidence Reuse
(SAVER)}, which organizes finalized PBRP evidence as an authenticated
dependency DAG. Request, authorization, transition, and checkpoint
evidence shared by multiple redactions is represented once and reused
only under the same authenticated state context. Challenged requests
may selectively escalate to independent PQZK verification without
revalidating unaffected shared dependencies.

\medskip
\noindent\textbf{Step 1: Authenticated Redaction Audit Index.}
For each finalized redaction record $RR_i$, the audit service derives
\begin{equation}
\tau_i^A=\mathsf{PRF}_{K_A}(TID_i)
\end{equation}
and constructs
\begin{equation}
E_i^A=
(\tau_i^A,RID_i,BID_e,PID_i,e,v_b,v_b',
h_i^{PBRP},ts_i^{red},ptr_i^A)
\end{equation}
where
\begin{equation}
h_i^{PBRP}
=
H_V(\texttt{PBRP}\parallel\mathsf{enc}(PBRP_i))
\end{equation}
and $ptr_i^A$ references its supporting evidence. RAI is append-only
with respect to finalized redactions.

Entries are partitioned into $S_A$ authenticated shards:
\begin{equation}
sid_i^A=H_A(\tau_i^A)\bmod S_A
\end{equation}
whose roots form
\begin{equation}
R_{\mathrm{RAI}}
=
\mathsf{Merkle.Root}(R_1^A,\ldots,R_{S_A}^A)
\end{equation}
A monotonically increasing $v_A$ identifies each snapshot. For audit
epoch $e_A$, the service constructs
\begin{equation}
CP_{e_A}^A
=
H_V(
\texttt{RAI-CP}\parallel e_A\parallel
R_{\mathrm{RAI}}\parallel v_A\parallel ts_A)
\label{eq:rai_cp}
\end{equation}
which is anchored to the PBN. Here, $e_A$ is distinct from the
authorization epoch $e$.

\medskip
\noindent\textbf{Step 2: PQ-Authenticated Audit Query Resolution.}
An authorized auditor $AU_j$ submits
\begin{equation}
Q_j^A=
(QID_j,qtype_j,\phi_j,scope_j,e_A,[t_s,t_e],n_j,ts_j)
\end{equation}
and
\begin{equation}
\sigma_j^A
\leftarrow
\mathsf{PQSIG.Sign}(sk_{AU_j},H(Q_j^A))
\end{equation}
After verifying the auditor, signature, freshness, and scope, the
service resolves
\begin{equation}
\mathcal{R}_{Q_j}
=
\{E_i^A:E_i^A\models Q_j^A\}
\end{equation}
against the selected RAI snapshot. Completeness is claimed only when
the index supports authenticated coverage of the requested predicate
and scope; otherwise, only the returned records are authenticated.

\medskip
\noindent\textbf{Step 3: Authenticated Verification-Dependency DAG.}
For the selected records, SAVER constructs
\begin{equation}
G_{Q_j}^{V}=(V_{Q_j},E_{Q_j})
\qquad
V_{Q_j}
=
V_{\mathrm{req}}\cup V_{\mathrm{auth}}
\cup V_{\mathrm{trans}}\cup V_{\mathrm{cp}}
\end{equation}
representing request validation, committee authorization, SFMR/PQCH
transitions, and finalized checkpoints, respectively.

Each node $v$ is content-addressed as
\begin{equation}
id(v)=
H_V(
\mathsf{type}(v)\parallel
\mathsf{enc}(\mathsf{payload}(v))\parallel
\mathsf{enc}(\mathsf{parents}(v)))
\label{eq:node_id}
\end{equation}
Request nodes contain $C_i^{VR}$, $\eta_i$, and
$PBRP_i^{auth}$; authorization nodes contain $Auth_e^B$ and
$\Gamma_e$; transition nodes contain $\Theta_b$ and $\Xi_b$; and
checkpoint nodes contain the corresponding $A_b'$ and RAI checkpoint.
Shared evidence therefore appears as a single dependency node.
The response carries each node once, and each record
references its request, authorization, transition, and checkpoint nodes
by $id(v)$; the auditor computes $id(v)$ once per node when the node is
received, not once per dependent record. Request nodes admitted in the
same VPS window share the window signature $\sigma_w^{V}$.

The service additionally constructs
\begin{equation}
MP_{Q_j}^A
\leftarrow
\mathsf{Merkle.MultiProof}
(\mathcal{R}_{Q_j},R_{\mathrm{RAI}})
\end{equation}
together with authenticated coverage evidence $\Psi_{Q_j}^A$ when
required and supported. The response is
\begin{equation}
Resp_j^A=(Body_j^A,\sigma_{\mathrm{Resp}}^A)
\end{equation}
where
\begin{equation}
\begin{aligned}
Body_j^A=\big(
&QID_j,H(Q_j^A),Ans_j,\mathcal{R}_{Q_j},MP_{Q_j}^A,\\
&G_{Q_j}^{V},\Psi_{Q_j}^{A},
R_{\mathrm{RAI}},v_A,CP_{e_A}^A\big)
\end{aligned}
\end{equation}
and
\begin{equation}
\sigma_{\mathrm{Resp}}^A
\leftarrow
\mathsf{PQSIG.Sign}
(sk_A,H(\mathsf{enc}(Body_j^A)))
\end{equation}

\medskip
\noindent\textbf{Step 4: State-Bound Verification Reuse.}
The auditor first verifies the response signature, query binding,
anchored RAI checkpoint, multiproof, and any required coverage
evidence, and then traverses $G_{Q_j}^{V}$ in dependency order.

Each node is associated with a typed state context
\begin{equation}
\kappa(v)=
\begin{cases}
(RID_i,h_i^{\mathrm{st}},e),
&v\in V_{\mathrm{req}}\\
(BID_e,\Gamma_e,e),
&v\in V_{\mathrm{auth}}\\
(b,v_b,v_b',\Theta_b),
&v\in V_{\mathrm{trans}}\\
(v_{\mathrm{RLI}}',v_A,CP_{e_A}^A),
&v\in V_{\mathrm{cp}}
\end{cases}
\end{equation}
and verification identifier
\begin{equation}
\chi(v)=
H_V(
\texttt{VERIFY-CTX}\parallel
id(v)\parallel\mathsf{enc}(\kappa(v)))
\label{eq:chi}
\end{equation}
After successful verification, the auditor records
\begin{equation}
\mathsf{VCS}[\chi(v)]=\mathsf{valid}
\end{equation}
where VCS denotes the \emph{Verified-Context Store}. Reuse is permitted
only when
\begin{equation}
\mathsf{Reuse}(v,\chi')
=1
\Longleftrightarrow
\chi'=\chi(v)
\land
\mathsf{VCS}[\chi(v)]=\mathsf{valid}
\label{eq:saver_reuse}
\end{equation}
Thus, any change to the evidence, dependencies, authorization context,
transition state, or checkpoint requires fresh verification.

For each distinct authorization node, the auditor verifies $C_e^B$,
$\Gamma_e$, and at least $t$ committee approvals. For each transition
node, it verifies
\begin{equation}
\mathcal{P}_b^{exec}\subseteq\mathcal{P}_b^{auth}
\end{equation}
the bindings to $\Theta_b$ and $\Xi_b$, and the finalized checkpoint.
Request-specific verification checks
\begin{equation}
I_i'=I_i\qquad
L_i=H(I_i\parallel D_i)\qquad
L_i'=H(I_i\parallel D_i')
\end{equation}
together with $\Pi_i^{SF}$ linking
$(pos_i,L_i,L_i')$ to the shared transition.
It also checks that $h_i^{\mathrm{snap}}$ is the anchored
checkpoint of batch $b$ at version $v_i^{\mathrm{snap}}$ with
$v_b-k\le v_i^{\mathrm{snap}}\le v_b$, and that
$h_i^{\mathrm{leaf}}=H(TID_i\parallel I_i\parallel D_i\parallel pos_i)$
for the old leaf $L_i=H(I_i\parallel D_i)$ of the executed transition;
together these show that the target leaf was unchanged between validation
and execution. PQCH preservation, Eq.~\eqref{eq:pqch_cont}, is
verified once per distinct batch transition.
The $z_Q$ distinct checks are independent and are
evaluated jointly as one matrix product over the query's transitions; a
failing transition still rejects only the records that depend on it. At
each distinct checkpoint node, the auditor recomputes $h_b^{\mathrm{CP}}$
from $A_b'$ and compares it with the digest anchored on the PBN for
$(b,v_b')$ (Phase~5, Step~5), so the response carries no checkpoint
signature.

\medskip
\noindent\textbf{Step 5: Dependency-Selective Verification Escalation.}
Routine verification of request $i$ uses
\begin{equation}
\mathcal{D}_i^{\mathrm{routine}}
=
\{
v_i^{\mathrm{req}},
v_e^{\mathrm{auth}},
v_b^{\mathrm{trans}},
v_b^{\mathrm{cp}},
v_A^{\mathrm{cp}}
\}
\end{equation}
At the request node, the auditor verifies
$\mathsf{AttVerify}(\alpha_i)=1$
(Eq.~\eqref{eq:attverify}), checking each distinct window signature
once,
and confirms $h_i^{PQ}=H(\pi_i^{PQ})$. This authenticates the VPS
validation decision but does not independently establish PQZK validity.

For a challenged or forensic request, SAVER expands only the
corresponding request dependency:
\begin{equation}
\mathcal{D}_i^{\mathrm{deep}}
=
\mathcal{D}_i^{\mathrm{routine}}
\cup
\{x_i,\pi_i^{PQ}\}
\end{equation}
and independently verifies
\begin{equation}
\mathsf{PQZK.Verify}
(vk_{\mathrm{PQZK}},x_i,\pi_i^{PQ})=1
\qquad
H(\pi_i^{PQ})=h_i^{PQ}
\end{equation}
Unchanged shared dependencies remain reusable according to
Eq.~\eqref{eq:saver_reuse}.

\medskip
\noindent\textbf{Step 6: Query-Specific Audit Decision.}
Finally,
\begin{equation}
\mathsf{AuditVerify}
(PP,Q_j^A,Resp_j^A)
\rightarrow
(ans_j,\beta_j,\{\beta_{j,i}\}_{i\in\mathcal{R}_{Q_j}},\ell_j)
\end{equation}
where $\beta_{j,i}\in\{0,1\}$ is the decision for returned
record $i$, computed from $\mathcal{D}_i^{\mathrm{routine}}$ (or
$\mathcal{D}_i^{\mathrm{deep}}$) only, so an invalid record does not
invalidate records whose dependencies remain valid;
$\beta_j\in\{0,1\}$ indicates whether the verification conditions
required by $qtype_j$ are satisfied for the whole result
($\beta_j=1$ only if every $\beta_{j,i}=1$) and
$\ell_j\in\{\mathsf{routine},\mathsf{deep}\}$ records the applied
verification level. For exhaustive queries, $\beta_j=1$ additionally
requires authenticated coverage; otherwise, only the authenticity of
the returned records is certified.

For $q$ returned redactions spanning $a$ authorization nodes, $z$
transition nodes, and $c$ checkpoint contexts, with $d$ requests
requiring deep verification, the resulting verification cost is
\begin{equation}
\begin{aligned}
C_{\mathrm{SAVER}}
\approx{}&
qC_{\mathrm{req}}
+aC_{\mathrm{auth}}
+zC_{\mathrm{trans}}\\
&+cC_{\mathrm{cp}}
+dC_{\mathrm{PQZK}}
\end{aligned}
\end{equation}
where typically $a,z,c\ll q$ when redactions share authorization and
execution evidence. SAVER thereby converts the sharing established by
CORA and SFMR into state-bound verification reuse while retaining
request-level escalation when independent policy verification is
required.

    
    
    % =====================================================================
    

\section{Security Analysis}
\label{sec:security}

We analyze VeRedact-PQ against the threat model of Section~\ref{subsec:threat}.
Let $\mathcal{A}$ be a quantum-polynomial-time (QPT) adversary.
The analysis focuses on the security-critical redaction path; compromise
of the underlying permissioned consensus beyond its prescribed fault
threshold, VPS compromise, endpoint compromise, denial-of-service,
traffic-analysis leakage, and implementation-level side channels are
outside the primary model.

\subsection{Security Model and Assumptions}

VeRedact-PQ relies on the following assumptions.

\begin{itemize}

\item \textbf{A1---Post-Quantum Authentication.}
PQSIG is existentially unforgeable under quantum chosen-message
attacks (EUF-qCMA).

\item \textbf{A2---Private Policy Verification.}
PQZK satisfies quantum computational soundness and zero knowledge for
relation $\mathcal{R}_P$.

\item \textbf{A3---Threshold Redaction.}
PQCH provides collision resistance against QPT adversaries, and
threshold adaptation is realized as $\mathcal{F}_{\mathrm{TA}}$
(Phase~1, Step~5), which reveals only the collision randomness
$r_b'$ and requires at least $t$ shares. Fewer than $t$ contemporaneously valid
trapdoor shares are compromised within each resharing interval.
Previous shares are securely erased after proactive resharing.

\item \textbf{A4---Authenticated Commitments.}
The domain-separated hash functions are collision resistant at the
required post-quantum security level, the PRFs are quantum secure, and
the authenticated Merkle structures are binding.

\item \textbf{A5---Committee and Ledger Thresholds.}
Fewer than $t$ members of the active committee $\mathcal{C}_e$ are
compromised, and the PBN satisfies its prescribed consensus
fault-tolerance condition.

\item \textbf{A6---Validation Service.}
The VPS correctly authenticates requesters, evaluates public policy
conditions, and verifies PQZK proofs before issuing $\alpha_i$.
Routine auditing therefore authenticates the VPS decision rather than
independently reproducing it; deep auditing can independently verify
the committed PQZK proof. A corrupt VPS could admit an ineligible
request, delay eligible ones, or misstate admission times; the first is
detected by deep auditing of the escalated requests, while the latter
affect liveness and fairness but not the authorization guarantees of
Phases~4--6.

\end{itemize}

Under these assumptions, we establish request admission security,
threshold authorization, authorization-preserving aggregation,
authorized state-transition integrity, end-to-end accountability,
state-bound audit reuse, and the post-quantum security boundary.

\subsection{State-Bound Request Admission}

\begin{theorem}[Authenticated and State-Bound Admission]
An adversary cannot obtain a valid Phase~3 validation attestation for
an unauthorized requester, policy, operation, or authenticated
transaction state except with negligible probability.
\end{theorem}

\emph{Proof:} Admission passes four checks, and bypassing any of them
breaks one assumption.
\begin{itemize}[leftmargin=*,nosep]
\item \emph{Requester.} The VPS accepts $R_i$ only with a valid PQ
signature $\sigma_i^R$ under the registered key $pk_r$
(Eq.~\eqref{eq:req_sig}); signing a new request without $sk_r$ breaks
A1.
\item \emph{Policy and state.} $R_i$ is tied to the committed policy
$C_{P_i}$ (Eq.~\eqref{eq:policy_commit}) and to the digest $h_i^{st}$ of
its target leaf and snapshot (Eq.~\eqref{eq:hleaf_hsnap}); changing
either while keeping these hashes breaks A4.
\item \emph{Private eligibility.} The PQZK proof $\pi_i^{PQ}$ is verified
on a statement $x_i$ that contains $C_{P_i}$ and $h_i^{st}$
(Eq.~\eqref{eq:pqzk_stmt}); an accepting proof without a valid witness
breaks A2.
\item \emph{Attestation.} The VPS signs the window root $R_w^V$ whose
leaf $\nu_i$ contains $C_i^{VR}$ (Eq.~\eqref{eq:attverify}); accepting an
attestation for a request the VPS did not admit requires forging
$\sigma_w^V$ (A1) or a false Merkle membership proof (A4).
\end{itemize}
Under A6, the VPS issues the attestation only after these checks, so
unauthorized admission succeeds only with negligible probability.

\subsection{Threshold Authorization and CORA Preservation}

\begin{theorem}[Threshold Authorization Security]
A coalition controlling fewer than $t$ members of
$\mathcal{C}_e$ cannot independently authorize and execute a valid
redaction except with negligible probability.
\end{theorem}

\emph{Proof:} Both authorization and execution need $t$ members.
\begin{itemize}[leftmargin=*,nosep]
\item \emph{Authorization.} Phase~4 accepts a batch only with $t$ PQ
signatures on the same batch commitment $C_e^B$
(Eq.~\eqref{eq:auth_batch}); a coalition of $c<t$ members obtains the
missing approvals only by forging an honest member's signature (A1).
\item \emph{Execution.} An adaptation needs $t$ trapdoor shares combined
inside $\mathcal{F}_{\mathrm{TA}}$; fewer than $t$ contemporaneously
valid shares cannot produce one (A3), and resharing with secure erasure
prevents collecting shares across intervals.
\end{itemize}

\begin{lemma}[CORA Authorization Preservation]
CORA cannot make an unvalidated request eligible for committee
authorization or enlarge the authority of a validated request.
\end{lemma}

\emph{Proof:} CORA only chooses which validated requests are processed
together.
\begin{itemize}[leftmargin=*,nosep]
\item \emph{Input.} It selects only from the validated queue
$\mathcal{Q}_e$, and every selected request must still pass the
leaf-freshness check $\mathsf{FreshLeaf}_i$ (Eq.~\eqref{eq:freshleaf})
before authorization; stale requests are excluded.
\item \emph{Binding.} The selected requests enter the authorization root
$R_e^{VR}$ through their leaves $\eta_i$ (Eq.~\eqref{eq:eta}), and their
positions enter the plan $\Gamma_e$ (Eq.~\eqref{eq:cora_plan}); both are
signed in $C_e^B$. Inserting, replacing, or relocating a request
afterwards changes $R_e^{VR}$ or $\Gamma_e$ and invalidates the
approvals (A1, A4).
\end{itemize}

\subsection{Authorized-Subset State Transition}

\begin{theorem}[SFMR Authorized-Subset Integrity]
Except with negligible probability, a finalized Phase~5 transition
cannot modify a position outside the committee-authorized execution
plan.
\end{theorem}

\emph{Proof:} Let $\mathcal{P}_b^{auth}$ be the positions of batch $b$
authorized by $\Gamma_e$ and $\mathcal{P}_b^{exec}$ those executed.
\begin{itemize}[leftmargin=*,nosep]
\item \emph{Positions.} Phase~5 executes only a subset of the authorized
positions (Eq.~\eqref{eq:subset}), so a stale request may be withheld
but no unauthorized position added.
\item \emph{Content.} RADC changes only $D_i$ and keeps the immutable
context $I_i$, so altering identifier, owner, type, policy, or timestamp
breaks A4.
\item \emph{Transition.} The SFMR commitment $\Theta_b$
(Eq.~\eqref{eq:theta}) binds the executed positions, both roots, and
$\Gamma_e$; substituting any of them breaks A4.
\item \emph{Adaptation.} Members contribute shares only after checking
the above, and the result must keep the chameleon hash unchanged
(Eq.~\eqref{eq:pqch_cont}); otherwise A3 or A5 is violated.
\item \emph{Checkpoint.} Each batch checkpoint digest is recovered from
the signed round root $R_e^{\mathrm{CP}}$ by its membership proof, so an
altered checkpoint breaks A1, A4, or A5.
\end{itemize}

\subsection{Freshness and Replay Resistance}

\begin{theorem}[State and Epoch Freshness]
Validation or authorization evidence generated for one authenticated
state or authorization epoch cannot be reused after its
target leaf has changed, with a snapshot that is not an anchored
checkpoint within the last $k$ versions of its batch, or for an expired
epoch except with negligible probability.
\end{theorem}

\emph{Proof:} The request is bound to its state digest, which is part
of the PQZK statement $x_i$ and of $C_i^{VR}$:
\begin{equation}
h_i^{st}=H(h_i^{leaf}\parallel h_i^{snap})
\end{equation}
so presenting a different leaf or snapshot part requires an $H$
collision (A4). Phases~3--5 then require
$\mathsf{FreshLeaf}_i=1$ (Eq.~\eqref{eq:freshleaf}):
\begin{itemize}[leftmargin=*,nosep]
\item \emph{Changed leaf.} If the leaf now holds $D_i^{cur}\neq D_i$,
matching
\begin{equation}
H(TID_i\parallel I_i\parallel D_i^{cur}\parallel pos_i)=h_i^{leaf}
\end{equation}
requires a collision, so the request is rejected. If a later redaction
restores $D_i$, the leaf again holds exactly the validated content.
\item \emph{Old snapshot.} A snapshot that was never finalized, or is
more than $k$ versions old, has no anchored checkpoint; producing one
requires forging the round signature (A1), a false membership proof
(A4), or breaking A5.
\item \emph{Expired epoch.} $\alpha_i$, $C_e^B$, and $\Gamma_e$ carry the
authorization epoch, which is checked before execution; rewriting them
requires a collision or a signature forgery.
\end{itemize}
Redactions of other leaves of batch $b$ affect neither condition.

\subsection{End-to-End Redaction Accountability}

\begin{theorem}[Policy-Bound Redaction Accountability]
Every finalized VeRedact-PQ redaction can be cryptographically linked
to its validated request, threshold authorization, authorized execution
plan, executed state transition, and finalized checkpoint.
\end{theorem}

\emph{Proof:} The final proof $PBRP_i$ (Eq.~\eqref{eq:pbrp}) links each
step:
\begin{itemize}[leftmargin=*,nosep]
\item \emph{Request to authorization:} $PBRP_i^{auth}$ proves that the
request leaf $\eta_i$ is in the committee-signed root $R_e^{VR}$, whose
batch also commits to the plan $\Gamma_e$.
\item \emph{Authorization to transition:} the SFMR commitment $\Theta_b$
and the execution evidence $\Xi_b$ (Eq.~\eqref{eq:xi}) bind the executed
transition, and $\Pi_i^{SF}$ links the request's old and new leaves to
it.
\item \emph{Transition to record:} the redaction record commits to
\begin{equation}
h_i^{PBRP}=H_V(\texttt{PBRP}\parallel\mathsf{enc}(PBRP_i))
\end{equation}
which RAI and its anchored checkpoint authenticate.
\end{itemize}
Altering any link without changing its commitment breaks A4; forging an
authorization or checkpoint breaks A1 or A5.

\subsection{SAVER State-Bound Verification Reuse}

\begin{theorem}[State-Bound Verification-Reuse Soundness]
Assume that $\mathsf{VCS}[\chi(v)]=\mathsf{valid}$ is recorded only
after successful verification of evidence node $v$ under context
$\chi(v)$. A valid result cannot be reused for different evidence or a
different authenticated state context except with negligible
probability.
\end{theorem}

\emph{Proof:} A node's identifier $id(v)$ hashes its type, payload, and
parents (Eq.~\eqref{eq:node_id}), and its verification context
$\chi(v)$ hashes $id(v)$ with its state context $\kappa(v)$
(Eq.~\eqref{eq:chi}). A result is reused only when the contexts match
exactly (Eq.~\eqref{eq:saver_reuse}). Any change to the payload,
dependencies, authorization, transition, version, or checkpoint changes
$id(v)$ or $\kappa(v)$, so reuse under the same $\chi(v)$ requires an
$H_V$ collision (A4).

\begin{corollary}[Dependency-Selective Escalation]
Independent verification of a challenged request does not require
reverification of unchanged shared dependencies.
\end{corollary}

Routine auditing checks only the VPS attestation $\alpha_i$; deep
auditing additionally verifies the PQZK proof of the challenged request
(Eq.~\eqref{eq:pqzk_verify}), while its unchanged shared dependencies
remain reusable by Theorem~6.

\subsection{Audit Integrity, Completeness, and Privacy}

\begin{theorem}[Authenticated and Privacy-Preserving Audit]
Under A1, A2, and A4, an authorized auditor can verify authenticated
redaction evidence without learning the private policy witness, while
modification, substitution, or stale reuse of authenticated audit
evidence is detectable except with negligible probability.
\end{theorem}

\emph{Proof:}
\begin{itemize}[leftmargin=*,nosep]
\item \emph{Privacy.} Private credentials occur only in the witness
$w_i$; by zero knowledge, $\pi_i^{PQ}$ reveals nothing beyond the
validity of $\mathcal{R}_P(x_i,w_i)$ (A2). Routine audits never process
$w_i$, and deep audits see the proof, not the witness.
\item \emph{Integrity.} RAI entries are addressed by
\begin{equation}
\tau_i^A=\mathsf{PRF}_{K_A}(TID_i)
\end{equation}
authenticated under $R_{\mathrm{RAI}}$, bound to the anchored checkpoint
$CP_{e_A}^A$ (Eq.~\eqref{eq:rai_cp}), and returned in a PQ-signed
response, so modifying an entry, proof, root, version, or checkpoint
breaks A1, A4, or the ledger assumption.
\end{itemize}
The privacy claim is limited: query and access patterns and metadata
revealed by authorized results are not hidden.

\medskip
\noindent\textbf{Remark on Audit Completeness.}
Result-set completeness is claimed only with authenticated coverage
evidence $\Psi_{Q_j}^{A}$; otherwise only the integrity and freshness of
the returned records are certified.

\subsection{Post-Quantum Security Boundary}

\begin{theorem}[Post-Quantum Redaction Security]
Under A1--A6 and the stated committee and PBN thresholds, a QPT
adversary cannot cause an unauthorized VeRedact-PQ redaction to be
accepted except with negligible probability.
\end{theorem}

\emph{Proof:} An accepted redaction must pass every link of the chain,
each covered by an earlier result:
\begin{itemize}[leftmargin=*,nosep]
\item PQ requester authentication and state-bound PQZK validation
(Theorem~1);
\item $t$-of-$n$ PQ committee authorization (Theorem~2, Lemma~1);
\item authorized SFMR transition and threshold PQCH adaptation
(Theorems~2 and~3) on a fresh state and epoch (Theorem~4);
\item authenticated PQ audit evidence (Theorems~5--7).
\end{itemize}
Bypassing any link breaks one of A1--A6. Native consensus credentials may
remain classical: compromising one yields no requester signature, VPS
attestation, or $t$ approvals or PQCH shares, although violating the PBN
consensus assumption itself would affect ledger integrity independently
of VeRedact-PQ.

    % =====================================================================
    
    \section{Evaluation}
    \label{sec:evaluation}
    
    This section compares VeRedact-PQ analytically and experimentally with
    Schemes~\cite{ref1} (decentralized threshold redaction),
    \cite{ref13} (auditable redaction), \cite{ref27} (robust threshold
    redaction), and~\cite{ref34} (policy-hiding redaction)
    (Table~\ref{tab:comparison}), and with the lattice-based batch
    authentication Scheme~\cite{ref39}, which enters only Experiment~2.
    Each baseline enters only the experiments whose stage its paper
    defines.
    % ---------------------------------------------------------------------

\subsection{Computation Cost Analysis}
\label{subsec:computation-cost}

This subsection compares the dominant computation costs of VeRedact-PQ
with the primary baselines across request validation, authorization,
redaction execution, and audit verification, with the notation of
Table~\ref{tab:cost-notation}, where
$f_b=|\mathcal{F}_b(\mathcal{P}_b^{exec})|$ and
$F_e=\sum_{b\in\Omega_e}f_b$. Table~\ref{tab:comp-cost} compares the
dominant costs; N/A denotes a stage not supported by the scheme.

\begin{table}[t]
\centering
\caption{Notation Used in the Cost Analysis}
\label{tab:cost-notation}
\scriptsize
\renewcommand{\arraystretch}{1.08}
\setlength{\tabcolsep}{3pt}
\begin{tabular}{ll}
\hline
\textbf{Notation} & \textbf{Description} \\
\hline
$m$ & Requests selected by CORA \\
$|\Omega_e|$ & Transaction batches affected by CORA selection \\
$N$ & Leaves per transaction batch \\
$N_I$, $N_A$ & SA-RLI and RAI entries \\
$n$, $t$ & Committee size and authorization threshold \\
$f_b$, $F_e$ & SFMR frontier size for batch $b$, total frontier size \\
$n_Q$ & Redaction records returned by an audit query \\
$a_Q$ & Distinct authorization dependencies in audit result \\
$z_Q$ & Distinct transition dependencies in audit result \\
$c_Q$ & Distinct checkpoint contexts in audit result \\
$d_Q$ & Requests escalated to deep PQZK verification \\
$W$ & VPS admission-window size \\
$w_e$, $w_Q$ & Distinct VPS windows among the $m$ selected / $n_Q$ returned requests \\
$c$ & Challenged blocks in~\cite{ref13}, $c\le n_Q$ \\
\hline
$T_H$, $T_{\mathrm{PRF}}$ & Hash and PRF evaluation \\
$T_S$, $T_V$ & PQ signature generation and verification \\
$T_{ZP}$, $T_{ZV}$ & PQZK proof generation and verification \\
$T_{CH}$ & PQCH hash evaluation \\
$T_{PA}$, $T_{CB}$ & PQCH partial adaptation, share combination \\
$T_V^{c}$ & Classical signature verification \\
$T_{AD}^{c}$, $T_{CV}^{c}$ & Classical CH adaptation, verification \\
$T_{PA}^{c}$, $T_{CB}^{c}$ & Classical CH partial adaptation, combination \\
$T_{Pol}$ & Policy-based CH authorization operation \\
$T_{Acc}$ & RSA-accumulator membership/update operation \\
$T_{Tag}$ & Identity-based RSA tag operation~\cite{ref13} \\
$T_E^{c}$ & Classical modular exponentiation \\
\hline
\end{tabular}
\end{table}


\begin{table*}[t]
\centering
\caption{Computation-Cost Comparison of the Redaction Lifecycle}
\label{tab:comp-cost}
\footnotesize
\renewcommand{\arraystretch}{1.25}
\renewcommand{\cellset}{\renewcommand{\arraystretch}{1.45}}
\setlength{\tabcolsep}{3.2pt}
\begin{adjustbox}{max width=\textwidth}
\begin{tabular}{lcccc}
\toprule
\textbf{Scheme} &
\makecell{\textbf{Request Validation}\\\textbf{(per request)}} &
\makecell{\textbf{Authorization}\\\textbf{(per $m$ requests)}} &
\makecell{\textbf{Redaction Execution}\\\textbf{(per $m$ requests)}} &
\makecell{\textbf{Audit Verification}\\\textbf{(per $n_Q$ records)}} \\
\midrule

Scheme~\cite{ref1}
& $T_{Acc}+2T_V^{c}$
& $m\,t(T_{Acc}+2T_V^{c})$
& \makecell{$m(tT_{PA}^{c}+T_{CB}^{c}+T_{CV}^{c}+T_{Acc})$\\
$+O(m\log N)T_H$}
& $n_Q(T_{CV}^{c}+T_{Acc})$ \\

Scheme~\cite{ref13}
& N/A
& N/A
& \makecell{$m(T_{AD}^{c}+T_{CV}^{c}+3T_{Tag}+T_{Acc})$\\
$+O(mN)T_H$}
& $O(c)T_E^{c}$ \\

Scheme~\cite{ref27}
& N/A
& N/A
& $m(tT_{PA}^{c}+T_{CB}^{c}+T_{CV}^{c}+T_V^{c})$
& N/A \\

Scheme~\cite{ref34}
& N/A
& $mT_{Pol}$
& $m(T_{AD}^{c}+2T_{CV}^{c})$
& N/A \\

\midrule
\textbf{VeRedact-PQ}
& \makecell{$T_V+T_{ZV}+T_S/W+T_{\mathrm{PRF}}$\\
$+O(\log N_I+\log N+\log W)T_H$}
& \makecell{$w_eT_V+t(T_S+T_V)$\\
$+O(m\log W)T_H$}
& \makecell{$|\Omega_e|(tT_{PA}+T_{CB}+2T_{CH})+T_S$\\
$+O(m+F_e+\Delta_I+|\Omega_e|)T_H$}
& \makecell{$(w_Q+a_Qt+1)T_V+2z_QT_{CH}+d_QT_{ZV}$\\
$+O(n_Q\log W+a_Q+z_Q$\\
$+c_Q+\log N_A)T_H$} \\

\bottomrule
\multicolumn{5}{p{0.96\textwidth}}{\footnotesize
Requester-side proof/signature generation ($T_{ZP}+T_S$) is excluded
from server-side request validation. $F_e$ is the total SFMR frontier
size and $\Delta_I$ denotes the authenticated SA-RLI update work
($O(m\log N_I)$ without shared index-update paths). Routine SAVER
verification has $d_Q=0$; deep auditing adds $d_QT_{ZV}$.
Authorization, transition, and checkpoint evidence is verified once
per distinct SAVER dependency. The additional $T_V$ in audit
verification authenticates the audit-service response. In~\cite{ref1},
a block can be redacted only once and approvals carry no signature;
in~\cite{ref27}, the threshold is enforced inside adaptation but no
separate request-validation or audit stage is defined; the audit
in~\cite{ref13} produces an aggregate decision for the challenged
blocks; one-off attribute-key issuance in~\cite{ref34} is excluded.
For VeRedact-PQ, the VPS signs once per admission window
of $W$ requests and the checkpoint authority once per round; a checkpoint
context costs one digest recomputation against its anchored value. The
$2z_QT_{CH}$ term is evaluated as one batched matrix product. Scheme~\cite{ref39} defines authentication only and is compared in
Experiment~2.}
\end{tabular}
\end{adjustbox}
\end{table*}

As Table~\ref{tab:comp-cost} shows, every baseline repeats its
chameleon-hash work per redaction. VeRedact-PQ incurs a higher admission cost because each request
undergoes state-bound PQZK verification once before CORA selection.
Subsequent authorization uses the compact VPS attestation, while the
$t$ committee approvals are shared across the selected set:
\begin{equation}
C_{\mathrm{auth}}
=
w_eT_V+t(T_S+T_V)+O(m\log W)T_H
\label{eq:cauth}
\end{equation}
Thus, the threshold component amortizes to
$t(T_S+T_V)/m$ per request, and the VPS attestations to
$w_eT_V/m$.

CORA and SFMR reduce execution cost by exploiting batch and Merkle-path
locality. SFMR processes shared authentication structure according to
the total frontier size $F_e$, while threshold PQCH adaptation is
performed once per affected batch:
\begin{equation}
C_{\mathrm{PQCH}}
=
|\Omega_e|(tT_{PA}+T_{CB}+2T_{CH})
\qquad |\Omega_e|\leq m
\end{equation}
Hence, clustered requests reduce both $|\Omega_e|$ and $F_e$ relative
to independent processing; since $T_{PA}\gg T_H$, the reduction of
$|\Omega_e|$ dominates the saving.

SAVER similarly avoids repeated audit verification by sharing
state-bound dependencies:
\begin{equation}
C_{\mathrm{SAVER}}
=
n_QC_{\mathrm{req}}
+a_QC_{\mathrm{auth}}
+z_QC_{\mathrm{trans}}
+c_QC_{\mathrm{cp}}
+d_QT_{ZV}
\label{eq:saver_cost}
\end{equation}
where typically $a_Q,z_Q,c_Q\ll n_Q$. Routine auditing has $d_Q=0$;
only challenged requests require independent PQZK verification.

VeRedact-PQ thus does not assume cheaper PQ primitives; it invokes them
less often.


\subsection{Communication and Storage Cost Analysis}
\label{subsec:communication-cost}

Let $|R|$, $|\sigma|$, $|\pi|$, and $|h|$ denote the sizes of a
redaction request, PQ signature, PQZK proof, and hash digest,
respectively. Let $|\sigma^{c}|$, $|r^{c}|$, $|ID|$, $|BH|$, and
$|\mathbb{Z}_N|$ denote a classical signature, classical
chameleon-hash randomness, node identity, block header, and RSA-group
element. We use the same workload notation as in
Section~\ref{subsec:computation-cost}. Table~\ref{tab:comm-cost}
compares the dominant communication and persistent-state costs;
redacted payloads common to all schemes are omitted.

\begin{table*}[t]
\centering
\caption{Communication and Storage Cost Comparison}
\label{tab:comm-cost}
\footnotesize
\renewcommand{\arraystretch}{1.25}
\renewcommand{\cellset}{\renewcommand{\arraystretch}{1.45}}
\setlength{\tabcolsep}{3.2pt}
\begin{adjustbox}{max width=\textwidth}
\begin{tabular}{lcccc}
\toprule
\textbf{Scheme} &
\makecell{\textbf{Request Submission}\\\textbf{(per request)}} &
\makecell{\textbf{Authorization Evidence}\\\textbf{(per $m$ requests)}} &
\makecell{\textbf{On-Chain Redaction State}\\\textbf{(per $m$ requests)}} &
\makecell{\textbf{Audit Response}\\\textbf{(per $n_Q$ records)}} \\
\midrule

Scheme~\cite{ref1}
& $|R|+|\sigma^{c}|$
& $mt|ID|$
& $m(O(1)|h|+|r^{c}|)$
& $n_Q(|BH|+|\mathbb{Z}_N|)$ \\

Scheme~\cite{ref13}
& $|R|$
& N/A
& $m(|BH|+3|\mathbb{Z}_N|)$
& $c|BH|+(c+10)|\mathbb{Z}_N|$ \\

Scheme~\cite{ref27}
& N/A
& N/A
& $m|r^{c}|$
& N/A \\

Scheme~\cite{ref34}
& N/A
& N/A
& $m|r^{c}|$
& N/A \\

\midrule
\textbf{VeRedact-PQ}
& $|R|+|\sigma|+|\pi|$
& \makecell{$t|\sigma|+O(m)|h|$\\
$+O(F_e)|h|$}
& $|\sigma|+|\Omega_e|\,O(\log|\Omega_e|)|h|$
& \makecell{$(w_Q+a_Qt)|\sigma|+d_Q|\pi|$\\
$+O(n_Q\log W+a_Q+z_Q$\\
$+c_Q+\log N_A)|h|$} \\

\bottomrule
\multicolumn{5}{p{0.96\textwidth}}{\footnotesize
For VeRedact-PQ, complete PQZK proofs, PBRP objects, and detailed SFMR
evidence may be retained in the off-ledger evidence store, while their
digests or authenticated roots are anchored on-chain. Routine SAVER
auditing has $d_Q=0$; deep auditing transfers the PQZK proofs only for
the $d_Q$ challenged requests. Redacted payloads common to all schemes
are omitted. The $|QR_N|$ elements in~\cite{ref13} are represented by
$|\mathbb{Z}_N|$. The on-chain state of VeRedact-PQ is one
round checkpoint signature plus, per affected batch, its checkpoint
fields and a membership proof under $R_e^{\mathrm{CP}}$.}
\end{tabular}
\end{adjustbox}
\end{table*}


Table~\ref{tab:comm-cost} shows that VeRedact-PQ sends larger
per-request objects, a PQ signature and a PQZK proof, but only once at
validation; later phases reference the request through compact
commitments and $\alpha_i$, and committee authorization adds $t|\sigma|$
once per CORA-selected set. On-chain transition state grows with
$|\Omega_e|\leq m$ rather than $m$: one checkpoint per affected batch,
one round signature, and a $\lceil\log_2|\Omega_e|\rceil$-hash membership
proof per batch, while PBRP, PQZK, and SFMR objects stay off-ledger,
bound by anchored digests. For audits, SAVER sends each shared
dependency once, so the dominant response size is
\begin{equation}
\begin{aligned}
S_{\mathrm{SAVER}}
={}&O(n_Q\log W)|h|
+w_Q|\sigma|
+a_Q\,O(t|\sigma|)
+z_Q\,O(|h|)\\
&+c_Q\,O(|h|)
+d_Q|\pi|
+S_{\mathrm{RAI}}
\end{aligned}
\label{eq:saver_comm}
\end{equation}
where $S_{\mathrm{RAI}}$ is the query-scoped RAI evidence and typically
$a_Q,z_Q,c_Q\ll n_Q$. VeRedact-PQ thus does not shrink individual
post-quantum objects; it limits their replication.

    % ---------------------------------------------------------------------

\subsection{Performance Analysis}
\label{subsec:performance-analysis}

\subsubsection*{\textbf{Experimental Setup}}
All results come from executing the implemented systems; no timing is
simulated or extrapolated. Each experiment ran all its systems on one
AWS EC2 m7i.4xlarge instance (16 vCPUs, Intel Xeon Platinum 8488C,
64\,GiB RAM, Ubuntu~24.04); in Experiment~1, VeRedact-PQ's requesters
generated their proofs on a separate c8g.48xlarge instance (192
Graviton4 cores). The PBN is a 7-validator Hyperledger Besu QBFT network
with 2-s blocks and a 10-ms inter-node delay (Linux \texttt{tc netem}).
PQSIG is ML-DSA-65 (\texttt{liboqs}); $H$, $H_1$, $H_2$, $H_A$, $H_V$,
and all Merkle trees use domain-separated SHA3-256, and the PRFs use
HMAC-SHA3-256. PQZK is a winterfell~0.13 STARK (128-bit prime field with
quadratic extension, BLAKE3-256 commitments, blowup factor 32, 26 FRI
queries, 16 grinding bits) with at least 128-bit conjectured
soundness. PQCH is an SIS-based chameleon hash with a
Micciancio--Peikert trapdoor~\cite{ref10,ref17}, made $t$-of-$n$ by
Shamir-sharing each member's trapdoor component; since our threshold
mode samples the Gaussian perturbation spherically, we do not claim its
statistical security. PQ signatures are anchored as calldata, so they
count in transaction size and gas, but the testbed does not run their
validator-side verification, so finalization latencies exclude it.

The ledger holds $10^5$ transactions (256\,B--4\,KB) in 391 batches of
$N=256$ leaves; arrivals are Poisson, and targets follow a Zipf
distribution with skew $s$. Unless otherwise specified, $n=7$, $t=5$,
$B_{\min}=8$, $B_{\max}=256$, $T_{\max}=500$\,ms, $k=8$, $W=256$,
$T_{\mathrm{vps}}=500$\,ms, $S=S_A=16$, and $s=0.8$.
Schemes~\cite{ref1}, \cite{ref13}, \cite{ref27}, \cite{ref34},
and~\cite{ref39} were reimplemented in the same environment with their
own constructions (classical at about 128-bit security, and
lattice-based at NIST level~3 for Scheme~\cite{ref39}); stages a scheme
does not define are reported as unsupported. The committee members of
VeRedact-PQ and Schemes~\cite{ref1} and~\cite{ref27} share one process
but are timed as separate parties: a step costs its slowest member, and
each message round adds 10\,ms of network time. In Experiment~1,
Schemes~\cite{ref1} and~\cite{ref27} keep up to 16 redactions in flight
while their messages travel. Experiments~2--5 use three runs of 10
samples per point (2 for Scheme~\cite{ref13} at $n_Q=10^4$ in
Experiment~3), reported as the mean of the per-run medians with a 95\%
Student's $t$ interval and the pooled p95; Experiment~1 is one 60-s run
per point, with intervals from six 10-s sub-windows (throughput) and by
bootstrap (p95).
    % ------------------------------------


\subsubsection*{\textbf{Experiment 1: End-to-End Redaction Performance}}

This experiment tests whether CORA rounds and early validation sustain
more valid redactions while keeping invalid requests out. VeRedact-PQ
and Schemes~\cite{ref1}, \cite{ref13}, \cite{ref27}, and~\cite{ref34}
receive the same trace at 25--1,000 requests/s, of which 20\% are
invalid: 10\% \emph{stale} (referring to a content commitment that an
earlier redaction replaced) and 10\% \emph{absent} (targeting a
nonexistent transaction). We measure finalized valid redactions per
second, invalid requests accepted, and p95 latency from arrival to
finalization.

\begin{figure}[t]
\centering
\includegraphics[width=0.8\linewidth]{exp1_redaction_throughput.png}
\caption{End-to-End Redaction Performance}
\label{fig:exp1_redaction}
\end{figure}

As shown in Fig.~\ref{fig:exp1_redaction}, VeRedact-PQ peaks at 500
requests/s and leads every baseline from 250 requests/s; beyond the
peak, its throughput declines but stays above every baseline, as proofs age in the VPS queue and are re-proved
after failing $\mathsf{FreshLeaf}_i$. Scheme~\cite{ref27} levels off
near the confirmation capacity of the Besu network, since it anchors
every redaction, Scheme~\cite{ref1} levels off lower, and
Schemes~\cite{ref13} and~\cite{ref34} stay at a few redactions per
second. Up to 100 requests/s, however, Scheme~\cite{ref27} matches
VeRedact-PQ within the confidence intervals. VeRedact-PQ accepted no
invalid request, since Phase~3 rejects them before any committee, PQCH,
or ledger work, whereas Schemes~\cite{ref27} and~\cite{ref1} executed
many of them, and Schemes~\cite{ref13} and~\cite{ref34} some at the
lowest rate. Its p95 latency stays at about one block period plus
$T_{\max}$ up to 250 requests/s, slightly above Scheme~\cite{ref27}'s at
25 and 50 requests/s; higher p95 values are censored lower bounds.
Locality also reduces PQCH work: adaptations per finalized redaction
fall as the skew $s$ increases, whereas every baseline adapts once per
redaction. These are throughputs of a Python implementation on one
host, and our requesters see each checkpoint about one block period
before it is anchored, which can only lower VeRedact-PQ's
revalidations.

    % ---------------------------------------------------------------------

\subsubsection*{\textbf{Experiment 2: Multi-Party Authorization Performance}}

This experiment tests whether amortizing committee approval over a
round makes authorization cost independent of the committee size. We
measure the complete cost per request (Phase~3 validation plus the
round's Phase~4 authorization divided by $m$, including message rounds)
for $n\in\{4,\dots,64\}$, $t=\lfloor2n/3\rfloor+1$, and $m=64$, against
Schemes~\cite{ref1} (distributed authorization), \cite{ref34}
(policy-based), and~\cite{ref39} (lattice-based aggregate
authentication without committee, policy, or freshness check).

\begin{figure}[t]
\centering
\includegraphics[width=0.8\linewidth]{exp2_authorization_latency.png}
\caption{Multi-Party Authorization Performance}
\label{fig:exp2_auth}
\end{figure}

As shown in Fig.~\ref{fig:exp2_auth}, VeRedact-PQ's cost stays nearly
constant across $n$, since it pays the
committee term $t(T_S+T_V)$ of Eq.~\eqref{eq:cauth} and one message
round once per round. Scheme~\cite{ref1} costs an order of magnitude
more, dominated by one message round per request, and
Scheme~\cite{ref34}'s cost grows with its policy size.
Scheme~\cite{ref39}, however, is cheaper at every $n$, as it lacks the
private-policy predicate, committee, and freshness check
(Table~\ref{tab:comparison}), and runs in Python, so an optimized
version would be faster still. On computation alone, VeRedact-PQ is
also above Scheme~\cite{ref1}, because Phase~3 verifies one PQZK proof
per request; its advantage over Scheme~\cite{ref1} therefore relies on
the 10-ms message round. Requester-side work, including VeRedact-PQ's
PQZK proving, is excluded for all schemes. Table~\ref{tab:exp2_m} gives
the computation cost versus the round size $m$.

\begin{table}[t]
\centering
\caption{Complete Authorization Cost per Request (ms, Computation
Only) versus Round Size $m$ at $n=7$}
\label{tab:exp2_m}
\scriptsize
\renewcommand{\arraystretch}{1.08}
\setlength{\tabcolsep}{4pt}
\begin{tabular}{rcccc}
\hline
$m$ & VeRedact-PQ & Scheme~\cite{ref1} & Scheme~\cite{ref34} & Scheme~\cite{ref39} \\
\hline
1    & 1.697 & 0.335 & 7.857 & 1.704 \\
8    & 0.930 & 0.343 & 7.817 & 0.730 \\
32   & 0.835 & 0.332 & 7.739 & 0.628 \\
64   & 0.818 & 0.343 & 7.741 & 0.619 \\
128  & 0.808 & 0.343 & 7.722 & 0.624 \\
256  & 0.790 & 0.342 & 7.730 & 0.625 \\
\hline
\end{tabular}
\end{table}

    % ---------------------------------------------------------------------
\subsubsection*{\textbf{Experiment 3: Audit and Verification Efficiency}}

This experiment tests whether SAVER lowers the routine audit cost
($d_Q=0$) as the result size $n_Q$ grows from 10 to $10^4$. We measure
response generation plus verification per query against
Schemes~\cite{ref1} and~\cite{ref13} and two variants:
\textbf{Per-Record}, which sends and verifies complete evidence for
every record, and \textbf{Shared-Evidence}, which sends shared evidence
once but re-verifies it for each record.

\begin{figure}[t]
\centering
\includegraphics[width=0.8\linewidth]{exp3_audit_efficiency.png}
\caption{Audit and Verification Efficiency}
\label{fig:exp3_audit}
\end{figure}

As shown in Fig.~\ref{fig:exp3_audit}, VeRedact-PQ is the cheapest at
every $n_Q$: below both variants and
Scheme~\cite{ref1}, and orders of magnitude below Scheme~\cite{ref13},
whose cost is spent building an RSA-3072 witness for every challenged
block. The saving grows with the number of records per authorization
batch, as more records share evidence. VeRedact-PQ loses to
Scheme~\cite{ref13} on verification time and response size, since
Scheme~\cite{ref13} returns one constant-size aggregate proof; that
single decision, however, shows neither who authorized a redaction,
under which policy, or whether the transition was correct, nor which
record is invalid (Experiment~4). Its response is also larger than
Scheme~\cite{ref1}'s, though far smaller than Per-Record's
(Table~\ref{tab:exp3_split}).

\begin{table}[t]
\centering
\caption{Audit Cost Terms at $n_Q=10^4$ (Routine, $d_Q=0$)}
\label{tab:exp3_split}
\scriptsize
\renewcommand{\arraystretch}{1.08}
\setlength{\tabcolsep}{3pt}
\begin{tabular}{lccccc}
\hline
 & VeRedact-PQ & Shared-Ev. & Per-Record & \cite{ref1} & \cite{ref13} \\
\hline
Generation (ms)   & 60.7 & 55.9 & 62.2 & 2{,}496 & $2.79\times10^6$ \\
Verification (ms) & 5{,}701 & 8{,}925 & 9{,}793 & 34{,}750 & 61.4 \\
Response (MB)     & 12.5 & 12.5 & 214.2 & 4.96 & 0.021 \\
Decisions         & $n_Q$ & $n_Q$ & $n_Q$ & $n_Q$ & 1 \\
\hline
\end{tabular}
\end{table}

    % ---------------------------------------------------------------------
 \subsubsection*{\textbf{Experiment 4: Selective Verification and Evidence Integrity}}

This experiment tests whether VeRedact-PQ identifies invalid records
faster than the baselines. In 1\%, 5\%, or 10\% of the returned
records ($n_Q\in\{10,10^2,10^3,10^4\}$), at the same positions for every
system, we inject \emph{modified} evidence (a committed field altered),
\emph{substituted} evidence (valid evidence from another redaction or
dependency context), or \emph{stale} evidence (a real earlier
checkpoint of the same batch replayed), plus two faults specific to
VeRedact-PQ: forged committee approvals, each invalidating every record
of its authorization set, and forged PQZK proofs. Routine verification checks
$\alpha_i$ and the shared dependencies; deep verification also checks
the PQZK proof, so forged proofs are evaluated at escalation fractions
$d_Q/n_Q$ from 0 to 100\%. The baselines are Scheme~\cite{ref1} and
Scheme~\cite{ref13}, whose Audit returns one decision per query; to
locate the invalid records, we follow a failed Audit with its per-block
Query, a combination rule of our own over two procedures of
Scheme~\cite{ref13}. We report verification time without reads of anchored values
from the PBN, invalid-record rejection, and valid-record acceptance.

\begin{figure}[t]
\centering
\includegraphics[width=0.8\linewidth]{exp4_verification_time.png}
\caption{Selective Verification and Evidence Integrity}
\label{fig:exp4_verify}
\end{figure}

As shown in Fig.~\ref{fig:exp4_verify}, with 1\% modified records,
VeRedact-PQ verifies 4.7--73$\times$ faster than Scheme~\cite{ref13}
and 3.9--6.4$\times$ faster than Scheme~\cite{ref1} at every $n_Q$. It
is the fastest in all 36 settings (three fault classes, three
fractions, four values of $n_Q$; at $n_Q=10$, 1\% and 5\% round to no
tampered record), by 4.4--81$\times$ and 3.3--6.6$\times$,
respectively, because SAVER verifies each shared
node once whereas Scheme~\cite{ref1} performs an RSA accumulator
exponentiation per record and Scheme~\cite{ref13} a Query per block.
VeRedact-PQ and Scheme~\cite{ref1} reject every invalid record and
accept every valid one; Scheme~\cite{ref13}'s Audit rejects every valid
record of a faulty query, and its Query, deciding per block, also
rejects valid records that share a block with an invalid one: 1.7\%,
12.7\%, and 47.1\% of them at 1\% faults ($n_Q=10^2$ to $10^4$), and
14.6\%, 48.1\%, and 88.9\% at 10\%. Forged
proofs are found in proportion to the escalation fraction, so
$d_Q/n_Q$ sets both the deep-verification cost and the detection rate.
Without faults, however, Scheme~\cite{ref13}'s Audit is the fastest for
$n_Q\geq10^2$, as its cost does not depend on $n_Q$, whereas
VeRedact-PQ's routine verification grows with $n_Q$
(Table~\ref{tab:exp4_nofault}), mostly in PQCH transition checks
(Table~\ref{tab:exp4_breakdown}). Reads of anchored values, excluded
from the figure, add to the times of VeRedact-PQ and
Scheme~\cite{ref1} but not of Scheme~\cite{ref13}, which reads none.

\begin{table}[t]
\centering
\caption{Verification Time Without Faults (ms)}
\label{tab:exp4_nofault}
\scriptsize
\renewcommand{\arraystretch}{1.08}
\setlength{\tabcolsep}{3pt}
\begin{tabular}{rccccc}
\hline
$n_Q$ & \multicolumn{2}{c}{VeRedact-PQ} & Scheme~\cite{ref1} & \multicolumn{2}{c}{Scheme~\cite{ref13}} \\
 & routine & deep & & Audit & Query/block \\
\hline
$10$   & 9.1       & 17.4       & 34.8       & 57.8 & 635.0      \\
$10^2$ & 62.5      & 122.9      & 339.3      & 58.4 & 4{,}216.1  \\
$10^3$ & 522.9     & 1{,}039.8  & 3{,}363.7  & 58.5 & 17{,}998.9 \\
$10^4$ & 5{,}262.3 & 10{,}556.0 & 33{,}557.6 & 60.2 & 25{,}313.3 \\
\hline
\end{tabular}

\smallskip
\parbox{0.95\linewidth}{\tiny Reads of anchored values excluded; mean of
the per-run medians over three runs. Without faults, Scheme~\cite{ref13}'s
Audit passes and no Query follows; the Query column runs it on every
block.}
\end{table}

\begin{table}[t]
\centering
\caption{Routine Verification-Time Breakdown at $n_Q=10^4$}
\label{tab:exp4_breakdown}
\scriptsize
\renewcommand{\arraystretch}{1.08}
\setlength{\tabcolsep}{4pt}
\begin{tabular}{lc}
\hline
Component & Time (ms) \\
\hline
RAI and response authentication & 108.9 \\
Request attestations & 328.1 \\
Committee approvals ($a_Q=$ 157) & 51.6 \\
Transition/PQCH ($z_Q=$ 7{,}568) & 4{,}763.2 \\
Checkpoints & 114.7 \\
\hline
of which reads of anchored values & 238.4 \\
Total & 5{,}507.5 \\
\hline
\end{tabular}

\smallskip
\parbox{0.95\linewidth}{\tiny No faults. Each component
is the mean of the per-run medians of its own timer, so the components
do not add up to the total, which also includes the per-record PBRP
digest and RAI-entry checks; reads of anchored values fall in the first
and fourth rows.}
\end{table}

    % ---------------------------------------------------------------------
 \subsubsection*{\textbf{Experiment 5: Permissioned Blockchain Finalization Overhead}}

This experiment measures the PBN anchoring load per finalized redaction
(transactions, bytes, and gas) under the workload of Experiment~1, and
the size, gas, and latency of VeRedact-PQ's finalization transaction for
$|\Omega_e|$ from 1 to 128 affected batches.

\begin{figure}[t]
\centering
\includegraphics[width=0.8\linewidth]{exp5_blockchain_overhead.png}
\caption{Permissioned Blockchain Finalization Overhead}
\label{fig:exp5_blockchain}
\end{figure}
 
\begin{table}[t]
\centering
\caption{VeRedact-PQ Finalization Transaction versus Affected Batches (7 Validators)}
\label{tab:exp5_size}
\scriptsize
\renewcommand{\arraystretch}{1.08}
\setlength{\tabcolsep}{4pt}
\begin{tabular}{rccc}
\hline
$|\Omega_e|$ & Size (B) & Gas & Finalization (s) \\
\hline
1   & 4{,}014  & 138{,}919   & 1.93 \\
8   & 6{,}254  & 383{,}901   & 1.83 \\
32  & 13{,}934 & 1{,}223{,}991 & 1.44 \\
64  & 24{,}174 & 2{,}343{,}793 & 0.93 \\
128 & 44{,}654 & 4{,}583{,}527 & 1.83 \\
\hline
\end{tabular}
\end{table}

As shown in Fig.~\ref{fig:exp5_blockchain}, each baseline anchors one
transaction per redaction, whereas VeRedact-PQ anchors two per round, so its transactions per redaction
fall as $2/m$ as rounds grow with the load. Its transactions carry
post-quantum objects and are larger, so when its rounds are small, it
anchors more bytes per redaction than Scheme~\cite{ref27} up to 100
requests/s and than Scheme~\cite{ref13} at 25 requests/s; from 250
requests/s, it anchors fewer than every baseline. Its gas is below every
baseline at every rate. The finalization transaction grows linearly
with the number of affected batches and finalized within one block
period (Table~\ref{tab:exp5_size}); since all validators share one
host, the testbed cannot resolve consensus overhead across validator
counts.



\section{Conclusion}
\label{sec:conclusion}

This paper presented VeRedact-PQ, a scalable post-quantum framework
for policy-bound redaction and verifiable auditing in permissioned
blockchains. VeRedact-PQ combines state-bound PQZK validation,
threshold committee authorization, and distributed PQCH adaptation
with RADC and authenticated indexing. CORA identifies sharing across
validated redaction requests, while SFMR exploits this locality through
shared Merkle processing and batch-level PQCH adaptation. PBRP binds
authorization to the resulting state transition, and SAVER extends
these dependencies to state-bound audit verification with selective
PQZK escalation. The evaluation demonstrates the trade-offs among post-quantum
validation, multi-party authorization, shared redaction execution,
audit verification, and permissioned-blockchain finalization.
VeRedact-PQ achieves scalability by reducing how often expensive
operations are repeated rather than assuming cheaper post-quantum
primitives. Future work will investigate deployment-scale workloads,
adaptive committee management, stronger privacy against access-pattern
leakage, and integration with additional post-quantum redaction
constructions.
    
 \section{Acknowledgement}
    This research has been supported by the Thammasat University Research Unit in Cyber Security.
    
     \begin{thebibliography}{99}
    
    \bibitem{ref1}
    C. Li, Q. Shen and Z. Wu, "Redactable Blockchain From Decentralized Chameleon Hash Functions, Revisited," in IEEE Transactions on Computers, vol. 74, no. 6, pp. 1911-1920, June 2025, doi: 10.1109/TC.2025.3544878.
    
    
    
    
    \bibitem{ref2}
    K. Huang, X. Zhang, Y. Mu, X. Wang, G. Yang, X. Du, F. Rezaeibagha, Q. Xia, and M. Guizani,
    ``Building Redactable Consortium Blockchain for Industrial Internet-of-Things,''
    \textit{IEEE Trans. Ind. Informat.}, vol. 15, no. 6, pp. 3670--3679, Jun. 2019,
    doi: 10.1109/TII.2019.2901011.
    
    \bibitem{ref3}
    W. Wang, L. Wang, J. Duan, X. Tong, and H. Peng,
    ``Redactable Blockchain Based on Decentralized Trapdoor Verifiable Delay Functions,''
    \textit{IEEE Trans. Inf. Forensics Security}, vol. 19, pp. 7492--7507, 2024,
    doi: 10.1109/TIFS.2024.3431917.
    
    \bibitem{ref4}
    H. Guo, W. Gan, M. Zhao, C. Zhang, T. Wu, L. Zhu, and J. Xue,
    ``PriChain: Efficient Privacy-Preserving Fine-Grained Redactable Blockchains in Decentralized Settings,''
    \textit{Chinese Journal of Electronics}, vol. 34, no. 1, pp. 82--97, 2025,
    doi: 10.23919/cje.2023.00.305.
    
    \bibitem{ref5}
    D. Zhang, J. Le, X. Lei, T. Xiang, and X. Liao,
    ``Secure Redactable Blockchain With Dynamic Support,''
    \textit{IEEE Trans. Dependable Secure Comput.}, vol. 21, no. 2, pp. 717--731, 2024,
    doi: 10.1109/TDSC.2023.3261343.
    
    \bibitem{ref6}
    X. Wu, X. Du, Q. Yang, N. Wang, and W. Wang,
    ``Redactable consortium blockchain based on verifiable distributed chameleon hash functions,''
    \textit{J. Parallel Distrib. Comput.}, vol. 183, Art. no. 104777, 2024,
    doi: 10.1016/j.jpdc.2023.104777.
    
    \bibitem{ref7}
    S. Aguincha, E. Nunes, S. Eisa, and M. L. Pardal,
    ``ChainGuards: Verification of Sensed Data using Permissioned Blockchain Technology,''
    arXiv preprint arXiv:2603.20769, 2026.
    
    \bibitem{ref8}
    D. Derler, K. Samelin, D. Slamanig, and C. Striecks,
    ``Fine-Grained and Controlled Rewriting in Blockchains: Chameleon-Hashing Gone Attribute-Based,''
    in \textit{Proc. 26th Annu. Netw. Distrib. Syst. Security Symp. (NDSS)}, 2019,
    doi: 10.14722/ndss.2019.23066.
    
    \bibitem{ref9}
    G. Ateniese, B. Magri, D. Venturi, and E. R. Andrade,
    ``Redactable Blockchain -- or -- Rewriting History in Bitcoin and Friends,''
    in \textit{Proc. 2nd IEEE Eur. Symp. Security Privacy (EuroS\&P)},
    Paris, France, pp. 111--126, 2017,
    doi: 10.1109/EuroSP.2017.37.
    
    \bibitem{ref10}
    C. Peng, H. Xu, and P. Li,
    ``Redactable Blockchain Using Lattice-based Chameleon Hash Function,''
    in \textit{Proc. 2022 Int. Conf. Blockchain Technology Information Security (ICBCTIS)},
    pp. 94--98, 2022,
    doi: 10.1109/ICBCTIS55569.2022.00032.
    
    \bibitem{ref11}
    C. Wu, L. Ke, and Y. Du,
    ``Quantum resistant key-exposure free chameleon hash and applications in redactable blockchain,''
    \textit{Information Sciences}, vol. 548, pp. 438--449, 2021,
    doi: 10.1016/j.ins.2020.10.008.
    
    \bibitem{ref12}
    K. Y. Chan, L. Chen, Y. Tian, and T. H. Yuen,
    ``Reconstructing Chameleon Hash: Full Security and the Multi-Party Setting,''
    in \textit{Proc. 19th ACM Asia Conf. Computer and Communications Security (AsiaCCS)},
    pp. 1076--1091, 2024,
    doi: 10.1145/3634737.3656291.
    
    \bibitem{ref13}
    X. Zhang, Z. Cai, K. Chen, G. Ha and C. Jia, "Efficient Auditing and Querying in Verifiable Redactable Blockchain: A Lightweight VDS Protocol with Integrity Verification," 2025 IEEE 24th International Conference on Trust, Security and Privacy in Computing and Communications (TrustCom), Guiyang, China, 2025, pp. 504-511, doi: 10.1109/Trustcom66490.2025.00062.
    
    \bibitem{ref14}
    K. Huang, X. Zhang, Y. Mu, F. Rezaeibagha, and X. Du,
    ``Scalable and redactable blockchain with update and anonymity,''
    \textit{Information Sciences}, vol. 546, pp. 25--41, 2021,
    doi: 10.1016/j.ins.2020.07.016.
    
    \bibitem{ref15}
    Y. Dong, Y. Li, Y. Cheng, and D. Yu,
    ``Redactable consortium blockchain with access control: Leveraging chameleon hash and multi-authority attribute-based encryption,''
    \textit{High-Confidence Computing}, vol. 4, no. 1, Art. no. 100168, 2024,
    doi: 10.1016/j.hcc.2023.100168.
    
    \bibitem{ref16}
    Y. Zhang, Z. Ma, S. Luo, and P. Duan,
    ``Dynamic Trust-Based Redactable Blockchain Supporting Update and Traceability,''
    \textit{IEEE Trans. Inf. Forensics Security}, vol. 19, pp. 821--834, 2024,
    doi: 10.1109/TIFS.2023.3326379.
    
    \bibitem{ref17}
    X. Wang, Y. Chen, X. Zhu, C. Li, and K. Fang,
    ``A Redactable Blockchain Scheme Supporting Quantum-Resistance and Trapdoor Updates,''
    \textit{Applied Sciences}, vol. 14, no. 2, Art. no. 832, 2024,
    doi: 10.3390/app14020832.
    
    \bibitem{ref18}
    J. B. Klamti and M. A. Hasan,
    ``Revocable policy-based chameleon hash using lattices,''
    \textit{Journal of Mathematical Cryptology}, vol. 18, no. 1, Art. no. 20230012,
    pp. 152--182, 2024,
    doi: 10.1515/jmc-2023-0012.
    
    \bibitem{ref19}
    X. Huang, Y. Wang, Y. Ding, Q. Wu, C. Yang, and H. Liang,
    ``Dynamically redactable Blockchain based on decentralized Chameleon hash,''
    \textit{Digital Communications and Networks}, vol. 11, no. 3, pp. 757--767, 2025,
    doi: 10.1016/j.dcan.2024.10.013.
    
    \bibitem{ref20}
    M. Miao, X. Yang, J. Wei, G. Tian, and W. Susilo,
    ``A Verifiable and Redactable Blockchain with Lightweight Storage and Permission Supervision,''
    \textit{Information}, vol. 17, no. 2, Art. no. 176, 2026,
    doi: 10.3390/info17020176.
    
    \bibitem{ref21}
    S. Fugkeaw, S. Sungchai, S. Nakprame, and P. Sreekongpan,
    ``Enabling Secure and Scalable GDPR-Compliant Blockchain-Based e-KYC With Efficient Redaction,''
    \textit{IEEE Access}, vol. 13, pp. 136834--136853, 2025,
    doi: 10.1109/ACCESS.2025.3594656.
    
    
    \bibitem{ref22}
    C. G. Harris, "A Redactable Blockchain Architecture for Regulatory Compliance," 2025 IEEE International Conference on Blockchain (Blockchain), Zhengzhou, China, 2025, pp. 1-6, doi: 10.1109/Blockchain67634.2025.00010.
    
    \bibitem{ref23}
    Z. Wu, L. Wang, X. Zhang and X. Feng, "EMT: Extended Merkle Tree Structure for Inserted Data Redaction in Permissioned Blockchain," in IEEE Transactions on Network Science and Engineering, vol. 12, no. 4, pp. 3025-3038, July-Aug. 2025, doi: 10.1109/TNSE.2025.3555979.
    
    \bibitem{ref24}
    J. W. Heo, G. Ramachandran and R. Jurdak, "Decentralised Redactable Blockchain: A Privacy-Preserving Approach to Addressing Identity Tracing Challenges," 2024 IEEE International Conference on Blockchain and Cryptocurrency (ICBC), Dublin, Ireland, 2024, pp. 215-219, doi: 10.1109/ICBC59979.2024.10634438.
    
    
    \bibitem{ref25}
    T. Sengupta, S. Chakraborty, and S. Sural,
    ``ReAcct: Redaction control for interoperable blockchains,''
    in \textit{Proc. 7th Conf. Blockchain Research Appl. Innovative
    Networks Services (BRAINS)}, Zurich, Switzerland, 2025,
    pp. 1--10, doi: 10.1109/BRAINS67003.2025.11302947.
    
    \bibitem{ref26}
    L. Yin, J. Shi, B. Xie and Y. Liu, "Rchain: A Universal and Efficient Redactable Blockchain Scheme," 2025 IEEE International Conference on Blockchain (Blockchain), Zhengzhou, China, 2025, pp. 235-240, doi: 10.1109/Blockchain67634.2025.00038.
    
    \bibitem{ref27}
    Z. Liu, B. Zhou and Y. Zhao, "Enhancing Redactable Blockchain With Robust and Efficient Threshold Redaction," in IEEE Transactions on Dependable and Secure Computing, vol. 23, no. 2, pp. 3238-3251, March-April 2026, doi: 10.1109/TDSC.2025.3634512.
    
    \bibitem{ref28}
    F. Wang, R. Dong, J. Cui, Q. Zhang and H. Zhong, "Blockchain-Based Anonymous and Accountable Data Redaction for the Industrial Internet of Things," in IEEE Transactions on Cloud Computing, doi: 10.1109/TCC.2026.3725253.
    
    \bibitem{ref29}
    S. Biswas, S. Chakraborty, and R. S. Chakraborty,
    ``Redacting without regret: A dual-layer framework for
    policy-compliant blockchains,''
    in \textit{Proc. 7th Conf. Blockchain Research Appl. Innovative
    Networks Services (BRAINS)}, Zurich, Switzerland, 2025,
    pp. 1--9, doi: 10.1109/BRAINS67003.2025.11302930.
    
    \bibitem{ref39}
    Q. Li, Y. Xia, X. Liu, X. Chen, Z. Liu, Z. Guo, Z. Liu, and L. Fang, ``LETA: A Lattice-Based Efficient and Traceable Privacy-Preserving Batch Authentication Scheme for Vehicle Platoon in VANETs,'' \textit{IEEE J. Sel. Areas Commun.}, vol. 44, 2026, doi: 10.1109/JSAC.2025.3636982.    
    \bibitem{ref31}
    Z. Yuqi, "Policy-Hiding Chameleon Hash with Traceability for Redactable Blockchains," 2025 22nd International Computer Conference on Wavelet Active Media Technology and Information Processing (ICCWAMTIP), Chengdu, China, 2025, pp. 1-4, doi: 10.1109/ICCWAMTIP68645.2025.11352637.
    
    \bibitem{ref32}
    S. Hu, M. Li, J. Weng, J. -N. Liu, J. Weng and Z. Li, "IvyRedaction: Enabling Atomic, Consistent and Accountable Cross-Chain Rewriting," in IEEE Transactions on Dependable and Secure Computing, vol. 21, no. 4, pp. 3883-3900, July-Aug. 2024, doi: 10.1109/TDSC.2023.3339675.
    
    \bibitem{ref33}
    L. Xue, H. Huang, F. Xiao, Q. Li and W. Wang, "A Controllable, Publicly Auditable, and Redactable Blockchain With a Main–Auxiliary Architecture," in IEEE Transactions on Dependable and Secure Computing, vol. 23, no. 2, pp. 1847-1864, March-April 2026, doi: 10.1109/TDSC.2025.3620854.
    
    \bibitem{ref34}
    J. Xue et al., "Attribute-Based Policy-Hiding Redactable Blockchain With Authorizable Verification for Energy Internet," in IEEE Internet of Things Journal, vol. 12, no. 15, pp. 29570-29583, 1 Aug.1, 2025, doi: 10.1109/JIOT.2025.3569699.
    
    
    \bibitem{ref35}
    K. Huang, X. Li, F. Rezaeibagha, L. Zhang and X. Zhang, "Time Updatable Policy-Based Chameleon Hash for Traceable and Accountable Redactable Blockchain," in IEEE Transactions on Information Forensics and Security, vol. 21, pp. 1470-1483, 2026, doi: 10.1109/TIFS.2026.3655919.
    
    \bibitem{ref36}
    J. Duan, W. Wang, L. Wang and L. Gu, "Controlled Redactable Blockchain Based on T-Times Chameleon Hash and Signature," in IEEE Transactions on Information Forensics and Security, vol. 19, pp. 7560-7572, 2024, doi: 10.1109/TIFS.2024.3436925. 
    
    \bibitem{ref37}
    W. -C. Kuo, D. -R. Lin and S. -X. Chen, "Identity-Based Chameleon Hashing from the SIS Assumption in a Central Authority Model," 2025 IEEE Conference on Dependable and Secure Computing (DSC), Taipei, Taiwan, 2025, pp. 1-6, doi: 10.1109/DSC65356.2025.11260874.
    
    \bibitem{ref38}
    M. Shirmohammadi, A. Islam and H. Karimipour, "An Intent-Based Networking Framework for Secure and Privacy-Compliant Machine Unlearning Using Meta-Learning and Redactable Blockchain," in IEEE Internet of Things Journal, vol. 13, no. 9, pp. 18167-18181, 1 May1, 2026, doi: 10.1109/JIOT.2025.3638966. 
    \end{thebibliography}
    
    \end{document}